\documentclass[a4paper,12pt,twoside]{book}
\raggedbottom 
\setlength{\headheight}{14.5pt}
\usepackage{pgfplots}
\pgfplotsset{compat=1.18}
\usepackage{subfiles}
\usepackage{subfiles}  
\usepackage{name}
\usepackage{amsthm}
\usepackage{amsmath}
\usetikzlibrary{matrix, positioning, calc}
\usepackage{mdframed}
\usepackage[most]{tcolorbox}
\usepackage{listings}
\usepackage[utf8]{inputenc}
\usepackage[T1]{fontenc}
\usepackage{xcolor}
\usepackage{listings}
\usepackage{tcolorbox}
\usetikzlibrary{positioning}
\usetikzlibrary{arrows.meta, shapes.geometric, positioning, decorations.pathreplacing, calc}
\usepackage[svgnames]{xcolor}   % per i colori come DodgerBlue, Crimson, ecc.
\usepackage{tikz}
\usetikzlibrary{positioning, calc}
\usepackage{listings}
% \usetikzlibrary{map}
\usepackage{booktabs}
\usetikzlibrary{shadings, decorations.pathmorphing}

\usepackage{tikz}
\usetikzlibrary{shapes.geometric, arrows, positioning}

\tikzstyle{startstop} = [rectangle, rounded corners, minimum width=3cm,
minimum height=1cm, text centered, draw=black]
\tikzstyle{decision} = [diamond, aspect=2, minimum width=3cm,
minimum height=1cm, text centered, draw=black]
\tikzstyle{proc} = [rectangle, rounded corners, minimum width=3.5cm,
minimum height=1cm, text centered, draw=black]
\tikzstyle{arrow} = [thick,->,>=stealth]


%  --- box per il codice --- 



% --- Definizione Colori ---
\definecolor{backcolour}{RGB}{250, 250, 250}
\definecolor{framegray}{RGB}{220, 220, 220}
\definecolor{codeblue}{RGB}{0, 0, 139}
\definecolor{codegreen}{RGB}{0, 100, 100}
\definecolor{codegray}{RGB}{128, 128, 128}
\definecolor{numgray}{RGB}{180, 180, 180}











\definecolor{doorwood}{HTML}{8B4513}
\definecolor{doorwood2}{HTML}{A0522D}
\definecolor{doorframe}{HTML}{5C4033}
\definecolor{doorknob}{HTML}{FFD700}
\definecolor{floorcolor}{HTML}{2C2C3A}
\definecolor{wallcolor}{HTML}{1A1A2E}
\definecolor{goatbg}{HTML}{2A1A1A}
\definecolor{carbg}{HTML}{1A2A1A}
\definecolor{goldtext}{HTML}{DAA520}
\definecolor{goatlight}{HTML}{F5F0E0}
\definecolor{goatmid}{HTML}{D4C8A8}
\definecolor{goatdark}{HTML}{8B7D5E}
\definecolor{goatshadow}{HTML}{5C5040}
\definecolor{horncolor}{HTML}{B8A070}
\definecolor{horndark}{HTML}{6B5D3E}
\definecolor{eyecolor}{HTML}{3B2510}
\definecolor{nosepink}{HTML}{D4A0A0}
\definecolor{hoofcolor}{HTML}{3D3020}
\definecolor{carred}{HTML}{CC2222}
\definecolor{carreddk}{HTML}{881515}
\definecolor{carredlt}{HTML}{EE4444}
\definecolor{cargray}{HTML}{555555}
\definecolor{carchrome}{HTML}{C0C0C0}
\definecolor{carwindow}{HTML}{88BBDD}
\definecolor{floatcol}{HTML}{4A90D9}
\definecolor{intcol}{HTML}{E85D4A}
\definecolor{lostcol}{HTML}{F5A623}
\definecolor{arrowcol}{HTML}{555555}
\definecolor{bgfloat}{HTML}{EAF0FA}
\definecolor{bgint}{HTML}{FDEAEA}







% =================== FINE COLORI =============================



% --- Stile Codice ---
\definecolor{funcpurple}{rgb}{0.5, 0, 0.5} % Definizione colore per le funzioni
\lstdefinestyle{exactstyle}{
    language=C++,
    backgroundcolor=\color{backcolour},
    commentstyle=\upshape\color{codegray},
    keywordstyle=\bfseries\color{codeblue},
    stringstyle=\color{codegreen},
    basicstyle=\ttfamily\small\upshape,  
    breakatwhitespace=false,
    breaklines=true,
    keepspaces=true,
    numbers=left,
    numbersep=15pt,
    numberstyle=\tiny\color{numgray},
    showspaces=false,
    showstringspaces=false,
    showtabs=false,
    tabsize=4,
    frame=none, 
    xleftmargin=25pt, 
    morekeywords=[2]{swap, main, printf, cout},
    keywordstyle={[2]\color{funcpurple}\bfseries},
}
\lstset{style=exactstyle}











\lstdefinestyle{stile_disegno}{
    frame            = tb,
    tabsize          = 4,
    numbers          = left,
    framesep         = 3pt,
    framerule        = 0.4pt,
    commentstyle     = \color{Green},
    keywordstyle     = \color{blue},
    stringstyle      = \color{DarkRed},
    backgroundcolor  = \color{WhiteSmoke},
    showstringspaces = false,
    numberstyle      = \small\color{Gray},
    xleftmargin      = 2em,
    framexleftmargin = 2em,
}











% --- Comando per codice inline ---
\newcommand{\inlinecode}[1]{\texttt{\textcolor{codeblue}{#1}}}









\definecolor{panna}{RGB}{218,208,200}

% Ambiente SOLO UNO, con nome diverso da altri già esistenti
\newtcbtheorem[
  number within=section
]{mydefinitiondu}{Formula}{%
  enhanced,
  colframe=defscol,
  boxrule=0.6pt,
  arc=3mm,              % angoli stondati (non usare 300mm)
  titlerule=0mm,
  toptitle=1mm,
  bottomtitle=1mm,
  fonttitle=\bfseries\large,
  coltitle=black,
  colbacktitle=white!30,
  colback=panna
}{defndu}

% Comandi di comodo
\NewDocumentCommand{\defndu}{m+m}{%
  \begin{mydefinitiondu}{#1}{}%
    #2%
  \end{mydefinitiondu}%
}

\NewDocumentCommand{\defnrd}{mm+m}{%
  \begin{mydefinitiondu}{#1}{#2}%
    #3%
  \end{mydefinitiondu}%
}











\newtcolorbox{strategiabox}{
  colback=white!3,
  colframe=black,
  boxrule=0.8pt,
  title=\textbf{Metodologia risolutiva proposta:},
  fonttitle=\bfseries,
  colbacktitle=white!15,
  coltitle=black
}

\newcommand{\strategia}[1]{%
  \begin{strategiabox}
    #1
  \end{strategiabox}
}


\newcommand{\cella}[2]{
    \node[draw, minimum width=1cm, minimum height=1cm, fill=#2, rounded corners=2pt] at (#1,0) {}
}
\newcommand{\num}[2]{
    \node[font=\large\bfseries] at (#1,0) {#2}
}

% Colori
\definecolor{ordinato}{RGB}{200,230,170}
\definecolor{scambio}{RGB}{255,220,150}
\definecolor{confronto}{RGB}{180,210,240}
\definecolor{neutro}{RGB}{230,230,230}















 
% Theorem System
% The following boxes are provided:
%   Definition:     \defn 
%   Theorem:        \thm 
%   Lemma:          \lem
%   Corollary:      \cor
%   Proposition:    \prop   
%   Claim:          \clm
%   Fact:           \fact
%   Proof:          \pf
%   Example:        \ex
%   Remark:         \rmk (sentence), \rmkb (block)
% Suffix
%   r:              Allow Theorem/Definition to be referenced, e.g. thmr
%   p:              Add a short proof block for Lemma, Corollary, Proposition or Claim, e.g. lemp
%                   For theorems, use \pf for proof blocks

% ======= Real examples : 

% \defn{Definition Name}{
%     A defintion.
% }

% \thmr{Theorem Name}{mybigthm}{
%     A theorem.
% }

% \lem{Lemma Name}{
%     A lemma.
% }

% \fact{
%     A fact.
% }

% \cor{
%     A corollary.
% }

% \prop{
%     A proposition.
% }

% \clmp{}{
%     A claim.
% }{
%     A reference to Theorem~\ref{thm:mybigthm}
% }

% \pf{
%    proof.

% \ex{
%     some examples. 
% }

% \rmk{
%     Some remark.
% }

% \rmkb{
%     Some more remark.
% }

\usepackage{xcolor}
\newtheorem{esempio}{Esempio}[chapter]
\newtheorem{osservazione}{Osservazione}[chapter]
















% Define custom colors
\definecolor{defscol}{HTML}{ecd8d7} %For definitions
\definecolor{asumscol}{HTML}{ecd8d7} %For Assumptions

\definecolor{rmkscol}{HTML}{313160} %For remarks
\definecolor{exmscol}{HTML}{e04b52} %For examples

\definecolor{lemscol}{HTML}{2c3943} %For Lemmes
\definecolor{thmscol}{HTML}{595765} %For Theorems
\definecolor{prpscol}{HTML}{9c98b1} %For proposition
\definecolor{corscol}{HTML}{dfd9fd} %For corrolaries

\definecolor{clmscol}{HTML}{165c58} %For claims
\definecolor{facscol}{HTML}{28a8a1} %For facts



% ============================
% Definition
% ============================
\newtcbtheorem[number within=section]{mydefinition}{Definizione}
{
    enhanced,
    frame hidden,
    titlerule=0mm,
    toptitle=1mm,
    bottomtitle=1mm,
    fonttitle=\bfseries\large,
    coltitle=black,
    colbacktitle=defscol!40!white,
    colback=defscol!20!white,
}{defn}

\NewDocumentCommand{\defn}{m+m}{
    \begin{mydefinition}{#1}{}
        #2
    \end{mydefinition}
}

\NewDocumentCommand{\defnr}{mm+m}{
    \begin{mydefinition}{#1}{#2}
        #3
    \end{mydefinition}
}

% Lista di definizioni raggruppate in un unico box
% Uso: \defnlist{Titolo}{ \item[Nome] Descrizione ... }
\NewDocumentCommand{\defnlist}{m+m}{
    \begin{mydefinition}{#1}{}
        \begin{itemize}[label=$\bullet$, leftmargin=1.5em]
            #2
        \end{itemize}
    \end{mydefinition}
}



\newenvironment{codeexample}[1][]{%
    \ifx\\#1\\%
        \begin{esempio}%
    \else%
        \begin{esempio}[#1]%
    \fi%
    \leavevmode\par
}{%
    \par\vspace{0.8\baselineskip}% spazio e a capo dopo l'esempio
    \end{esempio}%
}


\newenvironment{codebox}{%
    \begin{tcolorbox}[
        colback=backcolour,
        colframe=framegray,
        boxrule=0.5pt,
        sharp corners,
        left=5mm, right=5mm,
        top=5mm, bottom=5mm,
        width=\linewidth
    ]
}{%
    \end{tcolorbox}
}

\newtcbtheorem[number within=section]{mypropi}{Proprietà}
{
    enhanced,
    frame hidden,
    titlerule=0mm,
    toptitle=1mm,
    bottomtitle=1mm,
    fonttitle=\bfseries\large,
    coltitle=black,
     colbacktitle=blue!40!white,
    colback=blue!20!white,
}{propi}

\NewDocumentCommand{\propi}{m+m}{
    \begin{mypropi}{#1}{}
        \begin{itemize}
            #2
        \end{itemize}
    \end{mypropi}
}





% ============================
% Assumption
% ============================
\newtcbtheorem[use counter from=mydefinition]{myassumption}{Assunzione}
{
    enhanced,
    frame hidden,
    titlerule=0mm,
    toptitle=1mm,
    bottomtitle=1mm,
    fonttitle=\bfseries\large,
    coltitle=black,
    colbacktitle=asumscol!40!white,
    colback=asumscol!20!white,
}{asum}

\NewDocumentCommand{\asum}{m+m}{
    \begin{myassumption}{#1}{}
        #2
    \end{myassumption}
}

\NewDocumentCommand{\asumr}{mm+m}{
    \begin{myassumption}{#1}{#2}
        #3
    \end{myassumption}
}

% ============================
% Theorem
% ============================

\newtcbtheorem[number within=section]{mytheorem}{Teorema}{
    enhanced,
    frame hidden,
    titlerule=0mm,
    toptitle=1mm,
    bottomtitle=1mm,
    fonttitle=\bfseries\large,
    coltitle=black,
    colbacktitle=thmscol!40!white,
    colback=thmscol!20!white,
}{thm}

\NewDocumentCommand{\thm}{m+m}{
    \begin{mytheorem}{#1}{}
        #2
    \end{mytheorem}
}

\NewDocumentCommand{\thmr}{mm+m}{
    \begin{mytheorem}{#1}{#2}
        #3
    \end{mytheorem}
}

\newenvironment{thmpf}{
	{\noindent{\it \textbf{Dimostrazione.}}}
	\tcolorbox[blanker,breakable,left=5mm,parbox=false,
    before upper={\parindent15pt},
    after skip=10pt,
	borderline west={1mm}{0pt}{thmscol!40!white}]
}{
    \textcolor{thmscol!40!white}{\hbox{}\nobreak\hfill$\blacksquare$} 
    \endtcolorbox
}

\NewDocumentCommand{\thmp}{m+m+m}{
    \begin{mytheorem}{#1}{}
        #2
    \end{mytheorem}

    \begin{thmpf}
        #3
    \end{thmpf}
}

% ============================
% Lemma
% ============================

\newtcbtheorem[use counter from=mydefinition]{mylemma}{Lemma}
{
    enhanced,
    frame hidden,
    titlerule=0mm,
    toptitle=1mm,
    bottomtitle=1mm,
    fonttitle=\bfseries\large,
    coltitle=black,
    colbacktitle=lemscol!40!white,
    colback=lemscol!20!white,
}{lem}

\NewDocumentCommand{\lem}{m+m}{
    \begin{mylemma}{#1}{}
        #2
    \end{mylemma}
}

\newenvironment{lempf}{
	{\noindent{\it \textbf{Proof for Lemma}}}
	\tcolorbox[blanker,breakable,left=5mm,parbox=false,
    before upper={\parindent15pt},
    after skip=10pt,
	borderline west={1mm}{0pt}{lemscol!40!white}]
}{
    \textcolor{lemscol!40!white}{\hbox{}\nobreak\hfill$\blacksquare$} 
    \endtcolorbox
}

\NewDocumentCommand{\lemp}{m+m+m}{
    \begin{mylemma}{#1}{}
        #2
    \end{mylemma}

    \begin{lempf}
        #3
    \end{lempf}
}

% ============================
% Corollary
% ============================

\newtcbtheorem[use counter from=mydefinition]{mycorollary}{Corollario}
{
    enhanced,
    frame hidden,
    titlerule=0mm,
    toptitle=1mm,
    bottomtitle=1mm,
    fonttitle=\bfseries\large,
    coltitle=black,
    colbacktitle=corscol!40!white,
    colback=corscol!20!white,
}{cor}

\NewDocumentCommand{\cor}{+m}{
    \begin{mycorollary}{}{}
        #1
    \end{mycorollary}
}

\newenvironment{corpf}{
	{\noindent{\it \textbf{Proof for Corollary.}}}
	\tcolorbox[blanker,breakable,left=5mm,parbox=false,
    before upper={\parindent15pt},
    after skip=10pt,
	borderline west={1mm}{0pt}{corscol!40!white}]
}{
    \textcolor{corscol!40!white}{\hbox{}\nobreak\hfill$\blacksquare$} 
    \endtcolorbox
}

\NewDocumentCommand{\corp}{m+m+m}{
    \begin{mycorollary}{}{}
        #1
    \end{mycorollary}

    \begin{corpf}
        #2
    \end{corpf}
}

% ============================
% Proposition
% ============================

\newtcbtheorem[use counter from=mydefinition]{myproposition}{Proposizione}
{
    enhanced,
    frame hidden,
    titlerule=0mm,
    toptitle=1mm,
    bottomtitle=1mm,
    fonttitle=\bfseries\large,
    coltitle=black,
    colbacktitle=prpscol!30!white,
    colback=prpscol!20!white,
}{prop}

\NewDocumentCommand{\prop}{+m}{
    \begin{myproposition}{}{}
        #1
    \end{myproposition}
}

\newenvironment{proppf}{
	{\noindent{\it \textbf{Dimostrazione.}}}
	\tcolorbox[blanker,breakable,left=5mm,parbox=false,
    before upper={\parindent15pt},
    after skip=10pt,
	borderline west={1mm}{0pt}{prpscol!40!white}]
}{
    \textcolor{prpscol!40!white}{\hbox{}\nobreak\hfill$\blacksquare$} 
    \endtcolorbox
}



\NewDocumentCommand{\propp}{+m+m}{
    \begin{myproposition}{}{}
        #1
    \end{myproposition}

    \begin{proppf}
        #2
    \end{proppf}
}

% ============================
% Claim
% ============================

\newtcbtheorem[use counter from=mydefinition]{myclaim}{Claim}
{
    enhanced,
    frame hidden,
    titlerule=0mm,
    toptitle=1mm,
    bottomtitle=1mm,
    fonttitle=\bfseries\large,
    coltitle=black,
    colbacktitle=clmscol!40!white,
    colback=clmscol!20!white,
}{clm}

\NewDocumentCommand{\clm}{m+m}{
    \begin{myclaim*}{#1}{}
        #2
    \end{myclaim*}
}

\newenvironment{clmpf}{
	{\noindent{\it \textbf{Proof for Claim.}}}
	\tcolorbox[blanker,breakable,left=5mm,parbox=false,
    before upper={\parindent15pt},
    after skip=10pt,
	borderline west={1mm}{0pt}{clmscol!40!white}]
}{
    \textcolor{clmscol!40!white}{\hbox{}\nobreak\hfill$\blacksquare$} 
    \endtcolorbox
}

\NewDocumentCommand{\clmp}{m+m+m}{
    \begin{myclaim*}{#1}{}
        #2
    \end{myclaim*}

    \begin{clmpf}
        #3
    \end{clmpf}
}

% ============================
% Fact
% ============================

\newtcbtheorem[use counter from=mydefinition]{myfact}{Fact}
{
    enhanced,
    frame hidden,
    titlerule=0mm,
    toptitle=1mm,
    bottomtitle=1mm,
    fonttitle=\bfseries\large,
    coltitle=black,
    colbacktitle=facscol!40!white,
    colback=facscol!20!white,
}{fact}

\NewDocumentCommand{\fact}{+m}{
    \begin{myfact}{}{}
        #1
    \end{myfact}
}

% ============================
% Proof
% ============================

\NewDocumentCommand{\pf}{+m}{
    \begin{proof}
        [\noindent\textbf{Dimostrazione.}]
        #1
    \end{proof}
}

% ============================
% Example
% ============================


\newenvironment{myexample}{
    \tcolorbox[blanker,breakable,left=5mm,parbox=false,
    before upper={\parindent15pt},
    after skip=10pt,
	borderline west={1mm}{0pt}{clmscol!40!white}]
}{
    \textcolor{clmscol!40!white}{\hbox{}\nobreak\hfill$\blacksquare$} 
    \endtcolorbox
}

\NewDocumentCommand{\exm}{m+m}{
    \begin{myexample}
	{\noindent{\it \textbf{Dimostrazione: #1 }}}\\ 
        #2
    \end{myexample}
}


% ============================
% Remark
% ============================


\NewDocumentCommand{\rmk}{+m}{
    {\it \color{rmkscol!40!white}#1}
}

\newenvironment{remark}{
    \par
    \vspace{5pt}
        {\par\noindent{\textbf{Spiegazione.}}}
        \tcolorbox[blanker,breakable,left=5mm,
        before skip=10pt,after skip=10pt,
        borderline west={1mm}{0pt}{rmkscol!20!white}]
}{
        \endtcolorbox
    \vspace{5pt}
}

\NewDocumentCommand{\rmkb}{+m}{
    \begin{remark}
        #1
    \end{remark}
}

% comandi
\makeatletter
\newcommand{\impliesleft}{\mathrel{\mathpalette\implies@left\relax}}
\newcommand{\implies@left}[2]{%
  \reflectbox{$\m@th#1\implies$}%
}






\tikzset{
  root/.style   = {draw, circle, fill=blue!20},
  level1/.style = {draw, rectangle, fill=red!20},
  level2/.style = {draw, rounded corners, fill=green!20}
}



% % Old styles hh 

% %--------------------------------------------------------------------------
% % 		THEOREMES STYLE
% %--------------------------------------------------------------------------

% %-------		DEFINITION 		-------	
% \newcounter{defo}[chapter]
% \newenvironment{defi}[1]{\refstepcounter{defo} 
% \begin{tcolorbox}[colback=yellow!20!white,colframe=yellow!15!black,title= \textbf{Définition \thechapter \ $\blacklozenge$ \thedefo \ | #1}]}{\end{tcolorbox}}

% %-------		THEOREME 		-------	
% \newcounter{th}[chapter]
% \newenvironment{thm}[1]{\refstepcounter{th}
% \begin{tcolorbox}[colback=mycolor!10,colframe=mycolor!10!black!80,title=\textbf{ Théorème \thechapter \ $\blacklozenge$ \theth \ | #1}]}{\end{tcolorbox}}

% %-------		PROPOSITION 		-------	
% \newcounter{prop}[chapter]
% \newenvironment{propt}[1]{\refstepcounter{prop}
% \begin{tcolorbox}[colback=mycolor!5,colframe=mycolor!10!linkscolor!80 ,title=\textbf{ Proposition \thechapter \ $\blacklozenge$ \theprop \ | #1}]}{\end{tcolorbox}}

% %-------		COROLLAIRE		-------
% \newcounter{cor}[chapter]
% \newenvironment{corr}[1]{\refstepcounter{cor}
% \begin{tcolorbox}[colback=mycolor!2,colframe=mycolor!10!linkscolor!40 ,title= Corolaire \thechapter \ $\blacklozenge$ \thecor \ | #1]}{\end{tcolorbox}}

% %-------		LEMME			-------
% \newcounter{lem}[chapter]
% \newenvironment{lemme}[1]{\refstepcounter{lem}
% \begin{tcolorbox}[colback=mycolor!10!blue!2,colframe=mycolor!50!blue!30,title=\textbf{Lemme \thechapter \ $\blacklozenge$ \thelem \ | #1}]}{\end{tcolorbox}}

% %-------		METHODE 			-------
% \newcounter{met}[chapter]
% \newenvironment{meth}[1]{\refstepcounter{met}
% \begin{tcolorbox}
% [enhanced jigsaw,breakable,pad at break*=1mm,
%  colback=red!20!white,boxrule=0pt,frame hidden,
%  borderline west={1.5mm}{-2mm}{red}] \color{red}
% {\textbf{Méthode \thechapter \ $\blacklozenge$ \themet \ | #1} } \color{black} \\ } {\end{tcolorbox}}

% %-------		REMARQUE 			-------
% \newcommand{\NB}[1]{
% \ \\
% \begin{tabular}{p{0.05\textwidth}p{0.80\textwidth}}
% \hline
% \vspace{-0.1cm} \includegraphics[scale=0.03]{./system/IDEA.png} & 	#1\\
% \hline
% \end{tabular}
% \ 
% \newline \ \newline
%  }

% %-------		EXEMPLE  		-------
% \newenvironment{exm}{ \begin{tcolorbox}
% [enhanced jigsaw,breakable,pad at break*=1mm,
%  colback=cyan!20!white,boxrule=0pt,frame hidden,
%  borderline west={1.5mm}{-2mm}{cyan}] \color{cyan}
% {\textbf{Exemple} } \color{black} \\ } {\end{tcolorbox}}


 
% \usepackage{mystyle}  % Pacchetto personalizzato con tutto in italiano
 % Pacchetto personalizzato con tutto in italiano
\begin{document}


\begin{titlepage}
    \vspace*{\fill}
    \centering
    {\Huge \textbf{Pel}\par}
    \vspace{1.5cm}
    {\Large Appunti del corso\par}
    \vspace{1cm}
    \vspace*{\fill}
\end{titlepage}

\newpage

% ============================
% Prefazione
% ============================
\chapter*{Nota al lettore }
\addcontentsline{toc}{chapter}{Prefazione}

Questi appunti raccolgono i contenuti del corso  \textbf{Programmazione  e laboratorio modulo 1 }. 

% Il percorso parte dai concetti fondamentali --- variabili, tipi di dato, strutture di controllo --- e arriva progressivamente a trattare argomenti più avanzati come la gestione della memoria, i puntatori, le reference e gli algoritmi di ordinamento.

% L'approccio adottato è orientato alla \textit{comprensione} prima che alla sintassi: per ogni costrutto si cerca di motivarne l'esistenza, analizzarne il comportamento e discuterne le implicazioni. Dove possibile, i concetti vengono accompagnati da esempi, diagrammi e analisi passo per passo.

\vspace{0.5cm}
\noindent\textbf{Struttura degli appunti.} Il materiale è organizzato nei seguenti capitoli:

\begin{itemize}
    \item \textbf{Capitolo 1 -- Fondamenti:} variabili, tipi di dato, reference, strutture di iterazione (\inlinecode{for}, \inlinecode{while}, \inlinecode{do-while}), passaggio dei parametri e primi algoritmi.
    % \item \textbf{Capitolo 2 -- Algoritmi di ordinamento:} bubble sort, con analisi dell'invariante, ottimizzazioni e complessità.
\end{itemize}

\vspace{0.5cm}
\noindent\textbf{Convenzioni tipografiche.}
\begin{itemize}
    \item Il codice inline è scritto in \inlinecode{carattere monospazio}.
    \item I blocchi di codice sono numerati e evidenziati.
    \item Le definizioni, i teoremi e le osservazioni sono racchiusi in riquadri colorati per facilitarne l'individuazione.
\end{itemize}

\tableofcontents

% ============================
% Capitolo 1: Fondamenti
% ============================
\chapter{Fondamenti}

\intro{}

Introduciamo i concetti di base del linguaggio C++ da cui partirà tutto il resto del corso.

\section{Definizioni fondamentali}
\defnlist{Concetti fondamentali}{
    \item \textbf{Tipi di dati:} Un tipo di dato in c++  definisce l'insieme dei valori che può assumere e le operazioni che possono essere eseguite su di essi. 
    
    \item \textbf{Oggetto:} Un oggetto  è un area di memoria a cui è associato un tipo.
    \item \textbf{Variabile:} Una variabile è un oggetto a cui è associato un identificatore.
}
Un puntatore è una variabile che contiene l'indirizzo di un'area di memoria. Nell'esempio seguente, \inlinecode{a} è un puntatore a un intero allocato dinamicamente tramite \inlinecode{malloc()}:



\begin{codeexample}[Allocazione di memoria dinamica]
\begin{lstlisting}
    int *a = (int*) malloc(sizeof(int)); 
\end{lstlisting}
\end{codeexample}
\begin{center}

    \begin{tikzpicture}
        % Variabile puntatore
        \draw (0,0) rectangle (0.8,0.8);
        \fill (0.4,0.4) circle (2pt);
        \node[below] at (0.4,0) {\texttt{a}};
        
        % Freccia
        \draw[->, thick] (0.8,0.4) -- (3,0.4);
        
        % Blocco di memoria (array)
        \draw (3,0) rectangle (7,0.8);
    \end{tikzpicture}
\end{center}

La funzione \inlinecode{malloc()} (ereditata dal linguaggio C) alloca un blocco di memoria della dimensione specificata da \inlinecode{sizeof(int)} e restituisce un puntatore generico (\inlinecode{void*}), che viene convertito a \inlinecode{int*} tramite un cast esplicito. La variabile \inlinecode{a} non contiene direttamente un valore intero, ma l'indirizzo della cella di memoria allocata. Per accedere al valore puntato si usa l'operatore di dereferenziazione \inlinecode{*a}.

\subsection{Casting}
Il casting consiste nella conversione di un valore da un tipo all'altro. Nel linguaggio C++, esistono due tipologie di conversione: 

\begin{itemize}
    \item \textbf{Conversione implicita:} il compilatore converte da solo i tipi nelle espressioni \virgolette{miste}, promuovendo il tipo più piccolo verso quello più grande.

    vediamo un esempio: 


\pagebreak

\begin{codeexample}[Casting implicito]
\begin{codebox}
\begin{lstlisting}
    int a = 7; 
    int b = 3; 
    double c= a*1.0/b;  //  a e b vengono promossi a in double  => c = 2.333...
\end{lstlisting}
\end{codebox}

\end{codeexample}


    Per esempio \inlinecode{int} $\rightarrow$ \inlinecode{double}, in questo caso non abbiamo perdita di informazione.
    \item  \textbf{Conversione esplicita (cast):} In questo caso viene forzata la conversione da un tipo all'altro.  In C++ viene utilizzata la sintassi. \inlinecode{static\_cast\textless type\textgreater(x)}. 

        vediamo un esempio: 
        \begin{codeexample}
        [Casting esplicito]
        \begin{codebox}
    \begin{lstlisting}
    int total = 5;
    int count = 2;
    double average;

    // senza cast: promozione automatica
    average = total / count;      // 5 / 2 = 2 (int), poi 2 => 2.0

    // con cast: controlli tu la conversione
    average = static_cast<double>(total) / count; // 5.0 / 2 = 2.5
    \end{lstlisting}
    \end{codebox}

        \end{codeexample}


    
\end{itemize}


\pagebreak 

\subsection{Perdita di informazione:}
\defn{Narrowing}{
    La perdita di informazione si ha quando si converte da un tipo \virgolette{capiente} a un più piccolo.
}











\begin{center}
\begin{tikzpicture}[
    every node/.style={font=\sffamily},
    >=Stealth
]

% --- FLOAT BOX (sinistra) ---
\node[
    draw=floatcol, line width=1.6pt, rounded corners=8pt,
    fill=bgfloat,
    minimum width=4.2cm, minimum height=5.8cm,
    label={[font=\sffamily\large\bfseries, text=floatcol]above:{\texttt{float}}}
] (floatbox) at (0,0) {};

% Contenuto float: parte intera + decimale
\node[
    draw=floatcol!70, line width=1pt, rounded corners=4pt,
    fill=white,
    minimum width=3.4cm, minimum height=1.6cm,
] (intpart) at (0, 1.1) {};
\node[font=\sffamily\bfseries\large, text=floatcol!80] at (0, 1.45) {3};
\node[font=\sffamily\footnotesize, text=arrowcol] at (0, 0.7) {parte intera};

\node[
    draw=lostcol!80, line width=1pt, rounded corners=4pt,
    fill=lostcol!10,
    minimum width=3.4cm, minimum height=1.6cm,
] (decpart) at (0, -1.1) {};
\node[font=\sffamily\bfseries\large, text=lostcol] at (0, -0.75) {.14159};
\node[font=\sffamily\footnotesize, text=arrowcol] at (0, -1.5) {parte decimale};



% --- FRECCIA CONVERSIONE ---
\draw[->, line width=2.5pt, color=arrowcol, shorten >=4pt, shorten <=4pt]
    (2.8, 0) -- node[above, font=\sffamily\small\bfseries, text=arrowcol, yshift=2pt] {(int)}
    (5.5, 0);

% Label troncamento
\node[font=\sffamily\footnotesize\itshape, text=intcol!80, align=center] at (4.15, -0.55) {troncamento};

% --- INT BOX (destra) ---
\node[
    draw=intcol, line width=1.6pt, rounded corners=8pt,
    fill=bgint,
    minimum width=4.2cm, minimum height=5.8cm,
    label={[font=\sffamily\large\bfseries, text=intcol]above:{\texttt{int}}}
] (intbox) at (8.3, 0) {};

% Contenuto int: solo parte intera
\node[
    draw=intcol!70, line width=1pt, rounded corners=4pt,
    fill=white,
    minimum width=3.4cm, minimum height=1.6cm,
] (intval) at (8.3, 1.1) {};
\node[font=\sffamily\bfseries\large, text=intcol!80] at (8.3, 1.45) {3};
\node[font=\sffamily\footnotesize, text=arrowcol] at (8.3, 0.7) {parte intera};

% Zona vuota / persa
\node[
    draw=gray!40, line width=1pt, rounded corners=4pt,
    fill=gray!5, dashed,
    minimum width=3.4cm, minimum height=1.6cm,
] (lostval) at (8.3, -1.1) {};
\node[font=\sffamily\footnotesize\itshape, text=gray!60, align=center] at (8.3, -1.1) {non rappresentabile};

% Valore sotto il box int


% --- INDICATORE PERDITA ---
\draw[->, line width=1.4pt, color=lostcol, dashed, bend left=30, shorten >=6pt, shorten <=6pt]
    (decpart.east) to
    node[below, font=\sffamily\small\bfseries, text=lostcol, yshift=-10pt, xshift=5pt] {PERSO!}
    (lostval.west);

% --- BARRA INFERIORE: riepilogo ---
\node[
    draw=arrowcol!30, line width=0.8pt, rounded corners=6pt,
    fill=arrowcol!5,
    minimum height=1.1cm,
    text width=11.5cm,
    align=center,
    font=\sffamily\small
] at (4.15, -4.8) {
    \textbf{Narrowing conversion:} \texttt{float} (32 bit, con decimali)
    $\longrightarrow$ \texttt{int} (32 bit, solo interi)
    $\Rightarrow$ la parte frazionaria viene \textbf{troncata}
};

\end{tikzpicture}
\end{center}

Nello standard di C++ 11 è stata introdotta l'inizializzazione  con \inlinecode{\{\}}, che impedisce la narrowing conversion generando un errore in fase di compilazione. 
\begin{codeexample}[Utilizzo dell'inizializzazione che impedisce il narrowing]
\begin{codebox}

\begin{lstlisting}
double x = 3.7;
int a = x;  // narrowing implicito: concesso
int b{x};   // narrowing vietato: errore di compilazione 
\end{lstlisting}
\end{codebox}

\end{codeexample}
\begin{itemize}
    \item La riga con \inlinecode{=} permette la conversione da double a int.
    \item La riga con \inlinecode{\{\}}, secondo lo standard C++ 11, deve rifiutare il narrowing; quindi, il compilatore restituisce un errore.
\end{itemize}

\subsection{Left value Right value}
\defnlist{Left value e Right value}{
    \item \textbf{Left value:} è un espressione che identifica una locazione di memoria persistente, cioè un oggetto che ha un indirizzo a cui si può accedere anche dopo la valutazione dell'espressione. Il nome \virgolette{left value} deriva dal fatto che storicamente  è ciò che può stare  a sinistra di un operatore di assegnamento. 
    \item \textbf{Right value:} è un'espressione temporanea senza un indirizzo persistente, che può stare solo a destra.
}













































\vspace{2.5cm}

\begin{center}

\begin{minipage}{15.5cm}
\begin{lstlisting}[style=stile_disegno,language = C++, escapechar = !,
    basicstyle = \ttfamily\bfseries\large, linewidth = .705\linewidth]
int x = 10;!\tikz[remember picture] \node[] (line1end){};!
int y;!\tikz[remember picture] \node[] (line2end){};!
y = 1 + x;!\tikz[remember picture] \node[] (line3end){};!
\end{lstlisting}
\end{minipage}

\end{center}
\begin{tikzpicture}[remember picture, overlay,
    every edge/.append style = {->, thick, >=stealth,
                                DimGray, dashed, line width = 1pt},
    every node/.append style = {align = center, minimum height = 10pt,
                                font = \bfseries\small},
    text width = 3.2cm]

  % ===== LINEA 1: int x = 10; =====
  % Nodo per "x" (lvalue) – approssimato sopra la riga 1
  \node[fill = blue!15, rounded corners = 3pt, text width = 2.6cm,
        above left = 1.5cm and 1.0cm of line1end]
        (L1) {lvalue\\[-2pt]{\scriptsize\normalfont locazione di \texttt{x}}};

  % Nodo per "10" (rvalue)
  \node[fill = orange!20, rounded corners = 3pt, text width = 2.6cm,
        right = 1.8cm of L1]
        (R1) {rvalue\\[-2pt]{\scriptsize\normalfont letterale \texttt{10}}};

  % Frecce
  \draw (L1.south) edge ([xshift=-2.2cm, yshift=0.25cm]line1end.west);
  \draw (R1.south) edge ([xshift=-0.2cm, yshift=0.25cm]line1end.west);

  % ===== LINEA 2: int y; =====

  % ===== LINEA 3: y = 1 + x; =====
  \node[fill = blue!15, rounded corners = 3pt, text width = 2.6cm,
        below left = 1.4cm and 1.0cm of line3end]
        (L3) {lvalue\\[-2pt]{\scriptsize\normalfont \texttt{y} riceve il risultato}};

  \node[fill = orange!20, rounded corners = 3pt, text width = 3.0cm,
        right = 0.3cm of L3]
        (R3a) {rvalue\\[-2pt]{\scriptsize\normalfont espressione \texttt{1\,+\,x} (= 11)}};

  \node[fill = Goldenrod!25, rounded corners = 3pt, text width = 3.4cm,
        right = 0.8cm of R3a]
        (R3b) {\texttt{x} \`e un lvalue\\usato come rvalue\\[-2pt]{\scriptsize\normalfont si legge il valore di \texttt{x}}};

  \draw (L3.north) edge ([xshift=-2.8cm, yshift=-0.25cm]line3end.west);
  \draw (R3a.north) edge ([xshift=-1.5cm, yshift=-0.25cm]line3end.west);
  \draw (R3b.north) edge ([xshift=-0.2cm, yshift=-0.25cm]line3end.east);

\end{tikzpicture}

\vspace{2.5cm}





\section{Passaggio dei parametri}

Prima di addentrarci su come funzionano i passaggi per parametro in C++, andiamo a definire le \textbf{reference}. 


\defn{Reference}{
In C++ una reference è un alias per una variabile, ossia un nome alternativo che fa da riferimento all'area di memoria occupato dalla variabile originale. Una volta inizializzata, una reference non può essere riassegnata per riferirsi ad un altro oggetto.
}

\begin{center}
\begin{tikzpicture}[
    cell/.style = {draw, minimum width = 2.5cm, minimum height = 1cm, 
                   font = \ttfamily\bfseries\large, fill = blue!10},
    lbl/.style = {font = \bfseries\ttfamily\large},
    arrow/.style = {->, thick, >=stealth, DimGray, line width = 1pt},
  ]

  % === Cella di memoria ===
  \node[cell] (mem) at (0, 0) {42};
  \node[above = 0.1cm of mem, font=\footnotesize\ttfamily\color{Gray}] 
    {0x7ffc2a04};

  % === Etichette variabili ===
  \node[lbl] (x) at (-3.5, 0.6) {x};
  \node[lbl, color=DarkRed] (ref) at (-3.5, -0.6) {ref};

  % Frecce
  \draw[arrow] (x.east) -- (mem.west |- x);
  \draw[arrow, DarkRed] (ref.east) -- (mem.west |- ref);

  % Annotazione a destra
  \draw[decorate, decoration={brace, amplitude=6pt, mirror}, thick]
    ([xshift=0.3cm]mem.north east) -- ([xshift=0.3cm]mem.south east)
    node[midway, right = 10pt, font=\small, align=left] 
    {stessa area\\di memoria};

  % Tipo sotto le etichette
  \node[below = 0.05cm of x, font=\footnotesize\color{Gray}] {\texttt{int}};
  \node[below = 0.05cm of ref, font=\footnotesize\color{DarkRed!70}] {\texttt{int\&}};

\end{tikzpicture}
\end{center}



In C++ il passagio dei parametri può essere fatto per valore  e per riferimento.

\begin{codeexample}[Swap passagio per valore]
Nel passaggio per valore , alla funzione viene passata una \textbf{copia} del valore
dell'argomento effettivo. Le modifiche ai parametri formali all'interno della funzione non
influenzano le variabili originali nel chiamante.

\begin{codebox}
\begin{lstlisting}[language=C++]
void swap(int a, int b){
    int c = a;
    a = b;
    b= c;
    // Qui a e b sono effettivamente scambiati,
    // ma sono solo copie locali!
}

int main() {
    int x = 5;
    int y = 3;

    swap(x, y); // Passiamo i valori 5 e 3

    cout << x << " " << y;
    return 0;
}
\end{lstlisting}
\end{codebox}

L'output del programma sarà \texttt{5 3}: \texttt{x} e \texttt{y} non vengono scambiati
perché \texttt{a} e \texttt{b} sono solo copie locali che vengono distrutte alla fine
della funzione.
\end{codeexample}




\begin{codeexample}[Passaggio per copia di una \texttt{struct}]
Consideriamo una struttura \texttt{st} definita come segue:

\begin{codebox}
\begin{lstlisting}[language=C++]
struct st {
    int    a;
    int*   b;
    double c;
};
\end{lstlisting}
\end{codebox}

Quando si passa una \texttt{struct} per valore, il parametro formale \texttt{par}
riceve una \textbf{copia bit per bit} di tutti i campi del parametro effettivo.
Quanti byte vengono copiati dipende dalla composizione della struttura e viene
calcolato \textbf{staticamente} dal compilatore tramite \texttt{sizeof()}.

Nel caso di \textbf{allocazione statica} del campo \texttt{b} (ad esempio
\texttt{int b[20]}), \texttt{sizeof(st)} vale \texttt{21 int + 1 double} e
l'intera sequenza di byte viene copiata. In questo scenario non ci sono problemi:
la copia è completa e autonoma.

Il discorso cambia con l'\textbf{allocazione dinamica}. Se \texttt{b} è un
puntatore a un array allocato con \texttt{malloc}, \texttt{sizeof(st)} vale
\texttt{int + puntatore + double}: viene copiato solo il \emph{valore} del
puntatore, non l'array a cui esso punta. Si creano così due \textbf{puntatori
confluenti}: \texttt{x.b} e \texttt{par.b} puntano alla stessa area di memoria,
il che genera un \textbf{problema di ownership}: chi è responsabile di liberare
quella memoria? Una chiamata a \texttt{free(par.b)} all'interno di \texttt{foo}
lascerebbe \texttt{x.b} come puntatore \emph{dangling}.

\begin{codebox}
\begin{lstlisting}[language=C++]
struct st {
    int    a;
    int*   b;
    double c;
};

void foo(st par) {
    par.a    = 12;
    par.b[0] = 33;   // modifica l'array originale! (puntatori confluenti)
    par.c    = 3.14;
}

int main() {
    st x;
    x.b = (int*) malloc(sizeof(int) * 20);
    foo(x);
    cout << x.a;     // stampa 0: x.a non e' stato modificato
    return 0;
}
\end{lstlisting}
\end{codebox}

Si distinguono quindi due tipi di copia:
\begin{itemize}
    \item \textbf{Shallow copy} (copia superficiale): viene copiato solo il
          \emph{valore} del puntatore; i due puntatori condividono la stessa
          area di memoria.
    \item \textbf{Deep copy} (copia profonda): viene copiato anche il
          \emph{contenuto} dell'area puntata; le due strutture sono completamente
          indipendenti.
\end{itemize}
\end{codeexample}



\begin{codeexample}[\texttt{malloc} vs \texttt{new}]
In C++ esistono due modi per allocare memoria dinamicamente: la funzione
\texttt{malloc} (ereditata dal C) e l'operatore \texttt{new} (introdotto dal C++).
La differenza fondamentale è che \texttt{malloc} è una \textbf{funzione di libreria},
mentre \texttt{new} è un \textbf{operatore del linguaggio}.

Questa distinzione ha una conseguenza pratica importante: essendo \texttt{new} un
operatore, è \textbf{tipato}. \texttt{malloc} restituisce un \texttt{void*}, cioè un
puntatore generico che va esplicitamente castato al tipo desiderato; \texttt{new}
invece restituisce direttamente un puntatore al tipo specificato (ad esempio
\texttt{int*}), senza bisogno di cast.

\begin{codebox}
\begin{lstlisting}[language=C++]
// malloc: restituisce void*, richiede cast esplicito
int* p1 = (int*) malloc(sizeof(int) * 10);
free(p1);

// new: restituisce int*, nessun cast necessario
int* p2 = new int[10];
delete[] p2;
\end{lstlisting}
\end{codebox}

Alla \texttt{new} corrisponde l'operatore \texttt{delete} (o \texttt{delete[]} per
gli array), che libera la memoria allocata dinamicamente, analogamente a
\texttt{free} per \texttt{malloc}.

A differenza di linguaggi come Java o Python, il C++ \textbf{non ha un garbage
collector}. In quei linguaggi, il garbage collector è un programma che si attiva in
modo trasparente, controlla tutti gli oggetti allocati in memoria e verifica se
esistono ancora puntatori che vi fanno riferimento: se non ne esistono, libera
automaticamente quella memoria. In C++ la gestione della memoria è invece
completamente a carico del programmatore: ogni area allocata con \texttt{new} deve
essere esplicitamente liberata con \texttt{delete}, pena \textbf{memory leak}.
\end{codeexample}

\section{Uso preferenziale dei cicli }

In C++(e in generale), si usano tre tipi di cilico in base a quando conosci il numero di iterazioni e quando vuoi fare il testo della condizione.  Quando si  sceglie tra \inlinecode{while}, \inlinecode{do…while},\inlinecode{for} (e range-based for), la domanda vera è sempre: che tipo di ripetizione voglio esprimere?


\subsection{Uso preferenziale del for}
Il costrutto \inlinecode{for} è il più adatto quando il numero di iterazioni è collegato a un contatore che parte da un valore iniziale, viene aggiornato regolarmente e rimane all'interno di un certo intervallo.

La forma generale è: 















\vspace{2.3cm}

\begin{center}

\begin{minipage}{15.5cm}
\begin{lstlisting}[style=stile_disegno,language = C++, escapechar = !,
    basicstyle = \ttfamily\bfseries\large, linewidth = .705\linewidth]
for(int i = 0; i < size ; i++)!\tikz[remember picture] \node[] (line1end){};!
\end{lstlisting}
\end{minipage}

\end{center}
\begin{tikzpicture}[remember picture, overlay,
    every edge/.append style = {->, thick, >=stealth,
                                DimGray, dashed, line width = 1pt},
    every node/.append style = {align = center, minimum height = 10pt,
                                font = \bfseries\small},
    text width = 3.2cm]

  % Etichetta Inizializzazione
  \node[fill = blue!15, rounded corners = 3pt, text width = 3.2cm,
        above left = 1.5cm and 6.5cm of line1end]
        (INIT) {Inizializzazione\\[-2pt]{\scriptsize\normalfont \texttt{int i = 0}}};

  % Etichetta Condizione
  \node[fill = orange!20, rounded corners = 3pt, text width = 2.8cm,
        above left = 1.5cm and 3cm of line1end]
        (COND) {Condizione\\[-2pt]{\scriptsize\normalfont \texttt{i < size}}};

  % Etichetta Incremento
  \node[fill = green!20, rounded corners = 3pt, text width = 2.6cm,
        above left = 1.5cm and -0.5cm of line1end]
        (INCR) {Incremento\\[-2pt]{\scriptsize\normalfont \texttt{i++}}};

  % Frecce
  \draw (INIT.south) edge ([xshift=-7.5cm, yshift=0.25cm]line1end.west);
  \draw (COND.south) edge ([xshift=-4.5cm, yshift=0.25cm]line1end.west);
  \draw (INCR.south) edge ([xshift=-1cm, yshift=0.25cm]line1end.west);

\end{tikzpicture}


e ha il vantaggio di riassumere in una sola riga tutte le informazioni sul ciclo: dove parte, quando si ferma, come si aggiorna il contatore.



Il range‑based for (\inlinecode{for (auto x : v)}) è un caso particolare ancora più leggibile quando vogliamo dire “per ogni elemento della collezione v”; nasconde i dettagli sugli indici e mette l'attenzione sui dati.

In sintesi, se possiamo formulare il problema come “ripeti questa azione per tutti i valori di un contatore” o “per ogni elemento di una struttura”, il for (o il range‑based for) è la scelta preferenziale.





















\subsection{While}
Il \inlinecode{while} è il ciclo più generale e si usa quando la durata della ripetizione dipende da una condizione logica che non necessariamente è legata a un semplice contatore.

Il ciclo \inlinecode{while}  viene utilizzato nei seguenti casi: 
\begin{itemize}
    \item \textbf{ Sentinel‑controlled:} Qui il ciclo continua finché il valore letto non è uguale a un valore sentinella scelto apposta (per esempio -1), che non rappresenta un dato valido ma un segnale di “fine input”. È la situazione tipica quando non sappiamo quanti valori l’utente inserirà: leggiamo e processiamo ogni valore, e ci fermiamo quando arriva il sentinel.
    \item \textbf{Flag‑controlled:} n un ciclo flag‑controlled la condizione di continuazione dipende da una variabile booleana (il flag) che rappresenta lo stato del problema, per esempio “abbiamo già trovato quello che cercavamo?”. All’inizio il flag è impostato in modo da far entrare nel ciclo, e durante le iterazioni il programma aggiorna il flag; quando il flag assume il valore che significa “finito”, il ciclo termina. Questo schema è utile quando la fine della ripetizione dipende da un evento logico interno (trovare un elemento, soddisfare un vincolo, ecc.).

    \item \textbf{EOF-controlled:}Nel caso EOF‑controlled, la condizione del \inlinecode{while} è legata alla disponibilità di nuovi dati in uno stream di input (file o tastiera). Un pattern tipico è \inlinecode{while (in >> x)} oppure \inlinecode{while (cin >> x)}, che fa proseguire il ciclo finché l’operazione di lettura ha successo. Questo è il modo preferenziale per elaborare tutti i dati di un file senza conoscere in anticipo quanti record contiene.
\end{itemize}



La forma generale è: 





\vspace{2.3cm}
\begin{center}
\begin{minipage}{15.5cm}
\begin{lstlisting}[style=stile_disegno, language = C++,
    escapechar = !,
    basicstyle = \ttfamily\bfseries\large,
    linewidth = .705\linewidth]
while (condizione)!\tikz[remember picture] \node[] (wline1end){};!
        corpo;!\tikz[remember picture] \node[] (wline2end){};!
        i++;!\tikz[remember picture] \node[] (wline3end){};!
\end{lstlisting}
\end{minipage}
\end{center}

\begin{tikzpicture}[remember picture, overlay,
    every edge/.append style = {->, thick, >=stealth,
                                DimGray, dashed, line width = 1pt},
    every node/.append style = {align = center, minimum height = 10pt,
                                font = \bfseries\small},
    text width = 3.2cm]

  % Etichetta Condizione
  \node[fill = orange!20, rounded corners = 3pt, text width = 3.2cm,
        above left = 1cm and 2.5cm of wline1end]
        (WCOND) {Condizione\\[-2pt]{\scriptsize\normalfont testata prima di ogni iterazione}};

  % Etichetta Corpo
  \node[fill = blue!15, rounded corners = 3pt, text width = 3.0cm,
        below left = 1.5cm and 3.5cm of wline2end]
        (WBODY) {Corpo\\[-2pt]{\scriptsize\normalfont eseguito finché la condizione è vera}};

  % Etichetta Contatore
  \node[fill = green!20, rounded corners = 3pt, text width = 3.0cm,
        below left = 0.8cm and -2.5cm of wline3end]
        (WCNT) {Contatore\\[-2pt]{\scriptsize\normalfont aggiorna la variabile di controllo}};

  % Frecce
  \draw (WCOND.south) edge ([xshift=-3.0cm, yshift=0.25cm]wline1end.west);
  \draw (WBODY.north) edge ([xshift=-2.0cm, yshift=-0.25cm]wline2end.west);
  \draw (WCNT.north) edge ([xshift=-1cm, yshift=-0.25cm]wline3end.west);

\end{tikzpicture}

\vspace{2.5cm}

In generale il \inlinecode{while} è la scelta migliore quando vogliamo esprimere un processo che continua \virgolette{finché una certa situazione rimane vera}, dove la situazone può dipendere da input, stato interno o dalla fine dei dati. 






\section{Do while }
Il \inlinecode{do...while} si distingue perché la condizione viene controllata alla dine dell'iterazione, quindi il corpo del ciclo viene eseguito almeno una volta in ogni caso. 
È particolarmente adatto quando la logica naturale è “esegui questo blocco, poi decidi se ripeterlo”: è il caso tipico dei menu interattivi, dove bisogna mostrare le opzioni almeno una volta prima di chiedere se l'utente vuole continuare.

La forma generale del ciclo \texttt{do...while} è:

\vspace{2.3cm}
\begin{center}
\begin{minipage}{15.5cm}
\begin{lstlisting}[style=stile_disegno, language = C++,
    escapechar = !,
    basicstyle = \ttfamily\bfseries\large,
    linewidth = .705\linewidth]
do {!\tikz[remember picture] \node[] (dwline1end){};!
    corpo;!\tikz[remember picture] \node[] (dwline2end){};!
    i++;!\tikz[remember picture] \node[] (dwline3end){};!
} while (condizione);!\tikz[remember picture] \node[] (dwline4end){};!
\end{lstlisting}
\end{minipage}
\end{center}

\begin{tikzpicture}[remember picture, overlay,
    every edge/.append style = {->, thick, >=stealth,
                                DimGray, dashed, line width = 1pt},
    every node/.append style = {align = center, minimum height = 10pt,
                                font = \bfseries\small},
    text width = 3.2cm]

  % Etichetta Corpo
  \node[fill = blue!15, rounded corners = 3pt, text width = 3.0cm,
        above left = 0.8cm and 2.5cm of dwline2end]
        (DWBODY) {Corpo\\[-2pt]{\scriptsize\normalfont eseguito almeno una volta}};

  % Etichetta Contatore
  \node[fill = green!20, rounded corners = 3pt, text width = 3.0cm,
        below left = 0.8cm and 2.5cm of dwline3end]
        (DWCNT) {Contatore\\[-2pt]{\scriptsize\normalfont aggiorna la variabile di controllo}};

  % Etichetta Condizione
  \node[fill = orange!20, rounded corners = 3pt, text width = 3.2cm,
        below left = 0.8cm and 2.5cm of dwline4end]
        (DWCOND) {Condizione\\[-2pt]{\scriptsize\normalfont testata alla fine di ogni iterazione}};

  % Frecce
  \draw (DWBODY.south) edge ([xshift=-2.5cm, yshift=0.25cm]dwline2end.west);
  \draw (DWCNT.north) edge ([xshift=-2.5cm, yshift=-0.25cm]dwline3end.west);
  \draw (DWCOND.north) edge ([xshift=-3.5cm, yshift=-0.25cm]dwline4end.west);

\end{tikzpicture}

\vspace{2.5cm}


















\begin{table}[H]
\centering
\begin{tabular}{lll}
\toprule
Situazione                 & Ciclo   & Motivo                          \\
\midrule
Iterazioni da 1 a N        & \inlinecode{for}     & Contatore naturale              \\
Input finché sentinel      & \inlinecode{while}   & Valore sentinella per terminare \\
Menu ripetuto all'utente   & \inlinecode{do...while} & Corpo eseguito almeno una volta \\
Scansione array con indice & \inlinecode{for}     & Indice da 0 a size-1           \\
Condizione logica complessa& \inlinecode{while}   & Test separato dal setup         \\
\bottomrule
\end{tabular}

\medskip
\raggedright
\textit{Un valore sentinella (sentinel) è un valore speciale non valido come dato normale, usato solo per indicare la fine dell'input o del ciclo.}
\end{table}





\section{Programmazione strutturata}
In questo corso e in \virgolette{generale} è preferibile utilizzare il paradigma della \textbf{programmazione strutturata}, questo paradigma impone di scrivere i programmi usando re costrutti di controllo, senza ricorrere a salti incondizionati come \inlinecode{goto, break,continue}.


 I tre costrutti fondamentali:
 \begin{itemize}
     \item \textbf{Sequenza:} Le istruzioni vengono eseguite una dopo l'altra, nell'ordine in cui appaiono nel codice, Non c'è salto, non c'è scelta: si esegue A, poi B, poi C.
     \item \textbf{Selezione:} In base a una condizione booleana si sceglie quale ramo eseguire. In C++ sono l' \inlinecode{if}, \inlinecode{if\ldots else} e lo \inlinecode{switch}. Solo uno dei rami viene eseguito, poi il controllo riprende dopo la struttura. 
     \item \textbf{Ripetizione:} Un blocco di istruzioni viene eseguito ripetutamente finchè una condizione è vera. In C++ sono il \inlinecode{while}, il \inlinecode{do\ldots  while} e il \inlinecode{for}.
     
 \end{itemize}
\thm{ Böhm–Jacopini (1966)}{
Ogni algoritmo calcolabile può essere espresso usando soltanto i tre costrutti di sequenza, selezione e iterazione, senza mai ricorrere a salti incondizionati.
}





\chapter{Vettori e stringhe}
\intro{}


\section{Array}
\subsection{Definizione e idea intuitiva di array}
\defn{Array}{
Un array è una collezione di un numero fisso di elementi, tutti dello stesso tipo, memorizzati in celle di memoria contigue.

}
Dal punto di vista sintattico, la dichiarazione di un array in C++ segue la forma generale:
\begin{codebox}
\begin{lstlisting}
  dataType arrayName[intExp];
\end{lstlisting}
\end{codebox}

dove \inlinecode{dataType} indica il tipo degli elementi, \inlinecode{arrayName} è il nome attribuito all’array e \inlinecode{intExp} è un'espressione costante positiva che specifica il numero di componenti. È importante osservare che la dimensione così indicata rimane fissa per tutta la durata di vita dell’array: una volta dichiarato, non è possibile aumentare o diminuire il numero di elementi senza creare una nuova struttura e copiarvi i dati.



\subsection{ Limiti degli array e motivazione per strutture più astratte}

Gli array forniscono un meccanismo essenziale per rappresentare sequenze di valori omogenei in memoria contigua, con accesso diretto tramite indice. Tuttavia, proprio le caratteristiche che li rendono semplici ed efficienti ne evidenziano anche alcuni limiti strutturali, che diventano rilevanti nella progettazione di programmi complessi.

Un primo limite riguarda la rigidità della dimensione. La lunghezza di un array deve essere nota al momento della dichiarazione e rimane invariata per tutta la vita dell'array.

Ciò obbliga il programmatore a stimare in anticipo il numero massimo di elementi necessari: una scelta troppo prudente porta a spreco di memoria, mentre una stima troppo ottimistica può risultare insufficiente in casi reali d'uso. In altre parole, l’array è un contenitore statico, poco adatto a modellare collezioni la cui cardinalità può variare in modo significativo durante l'esecuzione.

Un secondo limite è legato alla mancanza di informazione incorporata sulla dimensione logica. Quando un array viene passato a una funzione, esso “decade” in un puntatore al suo primo elemento; la funzione riceve quindi solo un indirizzo di memoria, e non ha più accesso diretto alla dimensione dell’array originario. Di conseguenza, è necessario trasmettere esplicitamente un parametro aggiuntivo che rappresenti la lunghezza effettiva della sequenza, con il rischio di introdurre incoerenze tra i dati e le informazioni di supporto, oppure di incorrere in accessi fuori dai limiti qualora tali parametri non vengano gestiti con attenzione.


Inoltre, gli array non mettono a disposizione operazioni ad alto livello per la gestione dinamica dei dati. Non esistono primitive del linguaggio per inserire o rimuovere elementi all'interno della sequenza, per ridimensionare il contenitore o per verificarne automaticamente i limiti; tali operazioni devono essere simulate manualmente, spostando gli elementi e gestendo direttamente la memoria sottostante. Questa mancanza di astrazione aumenta la complessità del codice applicativo e rende più probabile l'introduzione di errori legati alla gestione della memoria, agli indici e alla coerenza dei dati.







Inoltre, gli array non mettono a disposizione operazioni ad alto livello per la gestione dinamica dei dati. Non esistono primitive del linguaggio per inserire o rimuovere elementi all’interno della sequenza, per ridimensionare il contenitore o per verificarne automaticamente i limiti; tali operazioni devono essere simulate manualmente, spostando gli elementi e gestendo direttamente la memoria sottostante. Questa mancanza di astrazione aumenta la complessità del codice applicativo e rende più probabile l’introduzione di errori legati alla gestione della memoria, agli indici e alla coerenza dei dati.



















\section{Vector}
\subsection{Idea generale}
\defn{Vector}{
 Un vector è un array dinamico: una sequenza ordinata di elementi omogenei memorizzati in celle di memoria contigue, la cui dimensione può variare durante l'esecuzione del programma. A differenza degli array statici, il \inlinecode{std::vector} integra la gestione automatica della memoria e mette a disposizione operazioni ad alto livello per inserire, rimuovere e accedere agli elementi.
}
 Dal punto di vista pratico , per utilizzare un \inlinecode{std::vector} in C++ la prima operazione necessaria è includere l'header della libreria standard che ne definisce l'interfaccia tramite la direttiva \inlinecode{\#include <vector>}.  




Una dichiarazione tipica è
\inlinecode{std::vector<int> v;}, che crea un vettore di interi vuoto,
la cui dimensione iniziale è 0.
In alternativa, la dichiarazione
\inlinecode{std::vector<int> v(100);}
costruisce un vettore con 100 elementi già presenti,
tutti inizializzati a 0.

Il punto chiave è che la dimensione logica del container
(non il tipo \inlinecode{std::vector<int>} in sé)
non è più una costante fissata staticamente come negli array,
ma uno stato interno dell'oggetto che può cambiare
durante l'esecuzione del programma.

\subsection{Struttura interna}

Dal punto di vista implementativo, \inlinecode{std::vector<T>} è un oggetto
che incapsula tre campi fondamentali:

\begin{itemize}
    \item un \textbf{puntatore} al buffer contiguo in memoria, dove sono
          effettivamente memorizzati gli elementi di tipo \inlinecode{T};
    \item \textbf{size}: un intero che rappresenta il numero di elementi
          attualmente presenti nel vettore;
    \item \textbf{capacity}: un intero che indica quanta memoria è stata
          riservata, ovvero quanti elementi possono essere contenuti prima
          che sia necessaria una riallocazione.
\end{itemize}

In generale vale sempre $\text{size} \leq \text{capacity}$:
quando si inserisce un elemento con \inlinecode{push\_back} e la size
raggiunge la capacity, il vector alloca automaticamente un nuovo buffer
più grande (tipicamente raddoppia la capacity), copia gli elementi

\subsection{Metodi principali}

Su questo oggetto \`e possibile chiamare i seguenti metodi:

\begin{itemize}
    \item \inlinecode{v.size()}: restituisce il numero di elementi logici
          attualmente contenuti nel vettore.

    \item \inlinecode{v.at(i)}: restituisce un riferimento all'elemento
          in posizione \inlinecode{i}, con controllo dei limiti.
          Se \inlinecode{i} \`e fuori dall'intervallo valido,
          lancia un'eccezione di tipo \inlinecode{std::out\_of\_range},
          causando un errore controllato anzich\`e un comportamento
          indefinito come con l'operatore \inlinecode{[]}.

    \item \inlinecode{v.resize(n)}: modifica la dimensione logica del
          vettore portandola a \inlinecode{n}. Se \inlinecode{n} \`e
          maggiore della size attuale, i nuovi elementi vengono
          inizializzati al valore di default del tipo (ad esempio
          \inlinecode{0} per \inlinecode{int}); se \inlinecode{n} \`e
          minore, gli elementi in eccesso vengono rimossi.

    \item \inlinecode{v.push\_back(x)}: aggiunge un nuovo elemento
          in coda al vettore, incrementando la size di 1.
          Se necessario, viene eseguita una riallocazione automatica
          del buffer interno.
\end{itemize}

\subsection{Passaggio dei parametri}



\begin{codeexample}[Passaggio dei parametri]
Un \inlinecode{std::vector} viene tipicamente passato alle funzioni
\textbf{per riferimento} , evitando la copia dell'intero
buffer e permettendo alla funzione di modificare il contenuto originale.



    \begin{codebox}

\begin{lstlisting}
void foo(vector<int>& v) {
    v.at(2) = 0;
}

int main() {
    vector<int> v {1, 2, 3, 4, 5};
    foo(v);
    cout << v.at(2);
    return 0;
}
\end{lstlisting}
\end{codebox}
\end{codeexample}

In questo esempio, \inlinecode{foo} riceve il vettore per riferimento
e modifica l'elemento in posizione 2. La chiamata
\inlinecode{cout<< v.at(2)} nel \inlinecode{main} stamper\`a quindi
\inlinecode{0}, confermando che la modifica ha effetto sull'oggetto
originale e non su una copia.

Il simbolo \inlinecode{\&} nella firma \inlinecode{vector<int>\& v} implica che
\inlinecode{v} \`e un \textbf{riferimento} al vettore originale: all'interno di
\inlinecode{foo} non viene creata alcuna copia, e qualunque modifica agli elementi
(ad esempio tramite \inlinecode{v.at(i)}) si riflette direttamente sul
vettore nel \inlinecode{main}.

Se invece si omettesse il simbolo \inlinecode{\&}, scrivendo
\inlinecode{void foo(vector<int> v)}, il parametro verrebbe passato
\textbf{per valore}: la funzione riceverebbe una copia locale del vettore,
tutte le modifiche rimarrebbero confinate a quella copia e
nessun effetto si propagherebbe al vettore originale nel chiamante.






























\section{Bubble sort}

Sia $v$ un vettore con una soglia arbitraria. Immaginiamo che la porzione del vettore che va da $0$ a soglia sia ordinata, mentre la porzione che va da soglia a $\text{size}-1$ sia disordinata.
\begin{center}
    \begin{tikzpicture}
        \draw[thick] (0,0) rectangle(10,1);
        \draw[thick](0,-0.4) --(0,-0.4) node [below]{$0$};
        \draw[thick](3,-0.4) --(3,-0.4) node [below]{soglia};
        \draw[thick](9,-0.4) --(9,-0.4) node [below]{$\text{size}-1$};
        \draw[very thick] (3,-0.4) -- (3,1.4);
    \end{tikzpicture}
\end{center}

Per ordinata intendiamo che gli elementi della porzione $[0, \text{soglia}]$ siano in ordine crescente e siano tutti pi\`u piccoli degli elementi che stanno nella porzione $[\text{soglia}, \text{size}-1]$, in cui gli elementi sono distribuiti in modo arbitrario. Inoltre, gli elementi in questa seconda porzione sono tutti maggiori o uguali agli elementi che stanno in $[0, \text{soglia}-1]$.

Potenzialmente ci troveremo nel caso in cui il vettore sia disordinato e quindi avremo la soglia uguale a zero. Questo ci fa leggere la condizione come la porzione che va da $0$ a $-1$, ovvero la porzione vuota, che consideriamo ordinata. La porzione $[0,\text{size}-1]$ \`e quindi disordinata, cio\`e l'intero vettore.

Pertanto, quando analizziamo il vettore per la prima volta, ci troviamo in questa situazione nella maggior parte dei casi. Quando invece il vettore \`e completamente ordinato, la soglia varr\`a \textit{size}, poich\'e la parte tra $0$ e $\text{size}-1$ sar\`a ordinata, e quella tra \textit{size} e $\text{size}-1$ sar\`a l'intervallo vuoto (visto che l'estremo superiore \`e minore dell'estremo inferiore).

Ad ogni iterazione del ciclo dovr\`o spingere la soglia verso destra.

\begin{center}
    \begin{tikzpicture}
        \draw[thick] (0,0) rectangle(10,1);
        \draw[thick](0,-0.4) --(0,-0.4) node [below]{$0$};
        \draw[thick](3,-0.4) --(3,-0.4) node [below]{soglia};
        \draw[thick](9,-0.4) --(9,-0.4) node [below]{$\text{size}-1$};
        \draw[very thick] (3,-0.4) -- (3,1.4);
        \draw[->] (3.5,1.5) -- (5.5,1.5);
    \end{tikzpicture}
\end{center}

Devo prendere l'elemento pi\`u piccolo che si trova a destra; devo portarlo nella posizione successiva alla soglia e, a questo punto, la soglia viene spostata a destra, poich\'e ho rispettato la regola: ho pescato sicuramente l'elemento che \`e maggiore o uguale a quelli nella porzione $\left[0, \text{soglia}-1\right]$ e l'ho posizionato subito dopo.

\begin{center}
    \begin{tikzpicture}
        \draw[thick] (0,0) rectangle(10,1);
        \draw[thick](0,-0.4) --(0,-0.4) node [below]{$0$};
        \draw[thick](3,-0.4) --(3,-0.4) node [below]{soglia};
        \draw[thick](9,-0.4) --(9,-0.4) node [below]{$\text{size}-1$};
        \draw[very thick] (3,-0.4) -- (3,1.4);
        \fill[gray!30] (2.99,0) rectangle (5,0.99);
        \draw[very thick] (5,-0.4) -- (5,1.4);
        \draw[thick](5.8,-0.4) --(5.8,-0.4) node [below]{new soglia};
        \draw[->] (3.5,1.5) -- (5.5,1.5);
    \end{tikzpicture}
\end{center}

\begin{codeexample}[Struttura base del bubble sort]
L'ordinamento viene implementato scorrendo il vettore dal fondo verso la soglia,
confrontando gli elementi a due a due e scambiandoli se necessario.

\begin{codebox}
\begin{lstlisting}[language=C++]
void bubble(vector<int> &v){
    for (int soglia = 0; soglia < v.size(); soglia++) {
        // La soglia deve prendere l'elemento
        // minore del range [soglia, size-1]
        // e metterlo in posizione soglia
    }
}
\end{lstlisting}
\end{codebox}
\end{codeexample}

L'ordinamento avviene se prendo il vettore e lo scorro nella direzione opposta rispetto a come scorre la soglia: dal fondo verso l'inizio. Posso considerare gli elementi a due a due: il penultimo con l'ultimo, e se l'ultimo \`e pi\`u piccolo, li scambio. Altrimenti, se sono gi\`a ordinati, non faccio nulla. Poi confronto il terzultimo con il penultimo, e cos\`i via.

Cosa succede quando arrivo all'elemento pi\`u piccolo del range soglia--size-1? L'elemento pi\`u piccolo verr\`a progressivamente spostato verso sinistra.

\begin{esempio}[Passata del bubble sort su un vettore]
Immaginiamo che la porzione che stiamo analizzando sia questa:
\begin{center}
    \begin{tikzpicture}
        \foreach \i/\val in {0/7,1/2,2/5,3/0,4/1} {
            \draw (\i,0) rectangle ++(1,1);
            \node at (\i+0.5,0.5) {\val};
        }
    \end{tikzpicture}
\end{center}

Comincio ad analizzare l'1 con lo 0:
\begin{itemize}
    \item Confronto l'1 con lo 0: sono disordinati.
    \begin{center}
        \begin{tikzpicture}
            \foreach \i/\val in {0/7,1/2,2/5,3/0,4/1} {
                \ifnum\i=3 \fill[blue!20] (\i,0) rectangle ++(1,1); \fi
                \ifnum\i=4 \fill[blue!20] (\i,0) rectangle ++(1,1); \fi
                \draw (\i,0) rectangle ++(1,1);
                \node at (\i+0.5,0.5) {\val};
            }
        \end{tikzpicture}
    \end{center}
    \item Confronto lo 0 con il 5: scambio.
    \item Confronto lo 0 con il 2: scambio.
    \item Confronto lo 0 con il 7: scambio.
\end{itemize}

Il valore 0 viene portato in prima posizione, e ottengo il seguente vettore:
\begin{center}
    \begin{tikzpicture}
        \foreach \i/\val in {0/0,1/7,2/2,3/5,4/1} {
            \draw (\i,0) rectangle ++(1,1);
            \node at (\i+0.5,0.5) {\val};
        }
    \end{tikzpicture}
\end{center}
\end{esempio}

\begin{codeexample}[Bubble sort completo]
Implementazione completa con la funzione \texttt{swap} e il doppio ciclo.

\begin{codebox}
\begin{lstlisting}[language=C++]
void swap(int &a, int &b){
    int c = a;
    a = b;
    b = c;
}

void bubble(vector<int> &v){
    for (int soglia = 0; soglia < v.size() - 1; soglia++) {
        for (int i = v.size() - 1; i > soglia; i--) {
            if (v.at(i) < v.at(i - 1))
                swap(v.at(i), v.at(i - 1));
        }
    }
}
\end{lstlisting}
\end{codebox}

\begin{itemize}
    \item \textbf{Dove devono fermarsi questi indici?}\\
    Partendo da \lstinline{size - 1}, il primo confronto \`e tra \lstinline{v[size-1]} e \lstinline{v[size-2]}, che \`e corretto. Se mi fermassi a \textit{soglia}, l'ultimo confronto sarebbe tra \textit{soglia} e \textit{soglia-1}. Tuttavia, l'invariante del ciclo ci garantisce che l'intervallo $[0, \text{soglia}-1]$ \`e ordinato. Quindi non serve confrontare gli elementi in questa parte. Concludiamo quindi che il ciclo pu\`o fermarsi a \lstinline{i > soglia}.

    \item \textbf{\`E corretto che \lstinline{i} parta da \lstinline{size - 1}?}\\
    S\`i, perch\'e se la soglia vale \lstinline{size - 1}, il ciclo interno non si esegue (nessun confronto), ed \`e corretto: un vettore con un solo elemento \`e gi\`a ordinato. Inoltre, se il vettore ha un solo elemento, l'intero ciclo esterno far\`a solo una iterazione inutile, che possiamo evitare usando la condizione \lstinline{soglia < v.size() - 1}.

    Mettere \texttt{size - 1} come condizione di arresto implica, secondo l'invariante del ciclo, che se la parte ordinata va da \texttt{0} a \texttt{size - 2} e la parte disordinata contiene un unico elemento, allora tutto l'array pu\`o essere considerato ordinato.
\end{itemize}
\end{codeexample}

\begin{tcolorbox}[
    colback=orange!10!white,
    colframe=orange!60!black,
    boxrule=0.5pt,
    sharp corners,
    left=5mm, right=5mm,
    top=3mm, bottom=3mm,
    title=\textbf{Attenzione a \texttt{size\_t}}
]
\lstinline{v.size()} restituisce un valore di tipo \lstinline{size_t}, che \`e \textit{unsigned}. Se il vettore \`e vuoto, \lstinline{v.size() - 1} non vale $-1$, ma causa un \textit{underflow} aritmetico, producendo un valore enorme ($2^{64}-1$ su architetture a 64 bit). Il confronto \lstinline{soglia < v.size() - 1} risulterebbe quindi sempre vero, portando a un comportamento indefinito. Per evitare il problema, \`e sufficiente aggiungere un controllo iniziale \lstinline{if (v.empty()) return;} oppure effettuare un cast esplicito \lstinline{soglia < (int)v.size() - 1}.
\end{tcolorbox}

\subsection*{Ottimizzazione: terminazione anticipata}

La versione base del bubble sort esegue sempre tutti i passaggi, anche se il vettore diventa ordinato prima di completarli. Possiamo ottimizzare l'algoritmo introducendo una variabile booleana \lstinline{check} che tenga traccia di un fatto semplice: \textit{durante l'ultimo passaggio, ho effettuato almeno uno scambio?}

L'idea \`e la seguente:
\begin{itemize}
    \item All'inizio di ogni passata, assumo che il segmento sia gi\`a ordinato (\lstinline{check = false}, cio\`e ``nessuno scambio effettuato'').
    \item Se durante la passata effettuo almeno uno scambio, imposto \lstinline{check = true}.
    \item Al termine della passata, se \lstinline{check} \`e rimasto \lstinline{false}, significa che non ho fatto alcuno scambio: il vettore \`e ordinato e posso terminare.
\end{itemize}

\textbf{Attenzione:} \`e fondamentale che l'assegnamento \lstinline{check = true} si trovi \textbf{dentro} il blocco \lstinline{if}. Se lo si mette fuori dall'\lstinline{if}, la variabile viene modificata a ogni iterazione indipendentemente dallo scambio, e l'ottimizzazione non funziona.

Poich\'e la variabile \lstinline{check} viene modificata nel corpo del ciclo e usata nella condizione di arresto, la scelta stilisticamente pi\`u corretta \`e usare un ciclo \lstinline{while} per il ciclo esterno:

\begin{codeexample}[Bubble sort ottimizzato con terminazione anticipata]
Uso di un ciclo \texttt{while} e della variabile \texttt{check} per terminare
anticipatamente quando il vettore \`e gi\`a ordinato.

\begin{codebox}
\begin{lstlisting}[language=C++]
void bubble(vector<int> &v) {
    if (v.empty()) return;

    bool check = true;
    int soglia = 0;

    while (soglia < (int)v.size() - 1 && check) {
        check = false;
        for (int i = v.size() - 1; i > soglia; i--) {
            if (v[i] < v[i - 1]) {
                swap(v[i], v[i - 1]);
                check = true;  // scambio avvenuto!
            }
        }
        soglia++;
    }
}
\end{lstlisting}
\end{codebox}
\end{codeexample}

\subsection*{Stabilit\`a}

Il bubble sort \`e un algoritmo \textbf{stabile}: se due elementi hanno lo stesso valore, il loro ordine relativo nel vettore originale viene preservato. Questo accade perch\'e lo scambio avviene solo quando \lstinline{v[i] < v[i-1]} (con disuguaglianza stretta): due elementi uguali non vengono mai scambiati tra loro.

\subsection*{Ordinamento in-place}

Il bubble sort ordina il vettore \textbf{in-place}, cio\`e senza richiedere memoria aggiuntiva proporzionale alla dimensione dell'input. L'unica memoria extra utilizzata \`e quella per le variabili ausiliarie (\lstinline{soglia}, \lstinline{i}, \lstinline{check} e la variabile temporanea dello \lstinline{swap}), quindi lo spazio aggiuntivo \`e $O(1)$.












\subsection{Il codice completo}
\section{String}
\inlinecode{std::string} \`e una classe della libreria standard che incapsula
una sequenza mutabile di caratteri, con molte istruzioni e operazioni disponibili.
Dal punto di vista strutturale \`e molto simile a un \inlinecode{vector<char>}.

\subsection{Differenze rispetto a \texttt{char*}}

A differenza delle C-string (\inlinecode{char*}), \inlinecode{std::string}
offre funzionalit\`a pi\`u ricche e sicure:

\begin{itemize}
    \item \textbf{Gestione automatica della memoria}: la stringa si ridimensiona
          automaticamente quando si aggiungono o rimuovono caratteri,
          senza bisogno di allocare o deallocare memoria manualmente.

    \item \textbf{Lunghezza in tempo costante}: \inlinecode{s.size()}
          restituisce il numero di caratteri in tempo $O(1)$, poich\'e
          la lunghezza \`e mantenuta come campo interno dell'oggetto.
\end{itemize}

\subsection{Input e output}

Per leggere e scrivere stringhe si utilizzano i flussi standard:
\inlinecode{cin} per l'input e \inlinecode{cout} per l'output.

\subsection{Operazioni fondamentali}

\subsubsection{Creazione}

\begin{itemize}
    \item \inlinecode{string s;} --- crea una stringa vuota.
    \item \inlinecode{string s = "ciao";} --- inizializza la stringa con il letterale \inlinecode{"ciao"}.
\end{itemize}

\subsubsection{Assegnamento}

\begin{itemize}
    \item \inlinecode{s = "hello";} --- sostituisce l'intero contenuto di \inlinecode{s} con \inlinecode{"hello"}.
    \item \inlinecode{s = s1;} --- copia il contenuto di \inlinecode{s1} in \inlinecode{s}.
\end{itemize}

\subsubsection{Lunghezza}

\inlinecode{s.size()} e \inlinecode{s.length()} sono equivalenti e restituiscono
il numero di caratteri della stringa in tempo $O(1)$.

\subsubsection{Accesso ai caratteri}

Come per il vector, \`e possibile accedere ai singoli caratteri tramite
\inlinecode{s.at(i)}, che effettua il controllo dei limiti, oppure
tramite l'operatore \inlinecode{s[i]}. Ad esempio:
\inlinecode{s.at(i) = 'z';} modifica il carattere in posizione \inlinecode{i}.

Un pattern comune \`e iterare sui caratteri con un range-for e modificarli
in place. L'esempio seguente converte tutti i caratteri minuscoli in maiuscoli:

\begin{codeexample}[Conversione minuscolo $\to$ maiuscolo]
    \begin{codebox}
        \begin{lstlisting}
for (char& c : s) {
    if (c >= 'a' && c <= 'z')
        c = c - 'a' + 'A';
}
        \end{lstlisting}
    \end{codebox}
\end{codeexample}

\subsubsection{Confronto e ordinamento}

Le stringhe supportano gli operatori di confronto standard:

\begin{itemize}
    \item \inlinecode{s1 == s2} --- vero se le due stringhe hanno la stessa
          lunghezza e gli stessi caratteri in ogni posizione.
    \item \inlinecode{s1 != s2} --- vero se almeno un carattere \`e diverso
          o le lunghezze sono diverse.
    \item \inlinecode{s1 < s2} --- confronto lessicografico: restituisce vero
          se \inlinecode{s1} precede \inlinecode{s2} nell'ordine del dizionario.
\end{itemize}

\subsection{Concatenazione e costruzione incrementale}

Per costruire una nuova stringa per concatenazione si usa l'operatore \inlinecode{+}:

\begin{itemize}
    \item \inlinecode{string s2 = s1 + s2;} --- crea una nuova stringa concatenando
          \inlinecode{s1} e \inlinecode{s2}.
\end{itemize}

Per la \textbf{concatenazione in-place} (senza creare oggetti temporanei) si usa
l'operatore \inlinecode{+=}:

\begin{itemize}
    \item \inlinecode{s1 += s2;} --- appende \inlinecode{s2} in fondo a \inlinecode{s1},
          modificandola direttamente.
    \item \inlinecode{s1 += '!';} --- aggiunge un singolo carattere in coda.
\end{itemize}

Per la \textbf{costruzione iterativa} carattere per carattere o blocco per blocco:

\begin{itemize}
    \item \inlinecode{s.push\_back('o');} --- aggiunge il carattere \inlinecode{'o'}
          in fondo alla stringa.
    \item \inlinecode{s.append("ciao");} --- appende la sottostringa
          \inlinecode{"ciao"} in fondo.
\end{itemize}

\subsection{Estrazione, ricerca e sostituzione}

\subsubsection{Operazioni di slicing con \texttt{substr}}

\begin{itemize}
    \item \inlinecode{s.substr(pos)} --- restituisce la sottostringa che va dalla
          posizione \inlinecode{pos} fino alla fine della stringa.
    \item \inlinecode{s.substr(pos, len)} --- restituisce la sottostringa di lunghezza
          \inlinecode{len} a partire dalla posizione \inlinecode{pos}.
\end{itemize}

\begin{codeexample}[Utilizzo di \lstinline{substr}]
    \begin{codebox}
        \begin{lstlisting}
string s = "ciao mondo";
string t = s.substr(5);    // "mondo"
string u = s.substr(0, 4); // "ciao"
        \end{lstlisting}
    \end{codebox}
\end{codeexample}

\subsubsection{Operazioni di ricerca con \texttt{find}}

\begin{itemize}
    \item \inlinecode{s.find("ciao")} --- restituisce l'indice della prima occorrenza
          della sottostringa \inlinecode{"ciao"} in \inlinecode{s}, oppure
          \inlinecode{string::npos} se non trovata.
    \item \inlinecode{s.find('a')} --- versione con singolo carattere.
\end{itemize}

L'esempio seguente estrae la seconda parola di una stringa cercando il primo spazio:

\begin{codeexample}[Prima parola dopo uno spazio]
    \begin{codebox}
        \begin{lstlisting}
size_t pos = s.find(' ');
if (pos != string::npos)
    string seconda = s.substr(pos + 1);
        \end{lstlisting}
    \end{codebox}
\end{codeexample}

\subsubsection{Sostituzione con \texttt{replace}}

\inlinecode{s.replace(pos, len, "nuovo\_testo")} sostituisce \inlinecode{len}
caratteri a partire dalla posizione \inlinecode{pos} con la stringa
\inlinecode{"nuovo\_testo"}.
In un'unica operazione cancella una porzione della stringa e inserisce
il testo di rimpiazzo, che pu\`o avere lunghezza diversa.

\subsubsection{Inserimento con \texttt{insert}}

\inlinecode{s.insert(pos, "xxx")} inserisce la stringa \inlinecode{"xxx"}
alla posizione \inlinecode{pos}, spostando in avanti i caratteri successivi.

\subsubsection{Cancellazione con \texttt{erase}}

\inlinecode{s.erase(pos, len)} cancella \inlinecode{len} caratteri
a partire dalla posizione \inlinecode{pos}.

L'esempio seguente mostra l'uso combinato di \inlinecode{insert}
e \inlinecode{erase} su una stringa:

\begin{codeexample}[Inserimento e cancellazione]
    \begin{codebox}
        \begin{lstlisting}
string s = "cane";
s.insert(1, "ial"); // s diventa "cialane"
s.erase(4, 3);      // s diventa "cial"
        \end{lstlisting}
    \end{codebox}
\end{codeexample}

\section{List}

\inlinecode{std::list<T>} \`e una \textbf{lista doppiamente concatenata}:
ogni nodo contiene un elemento e due puntatori, rispettivamente al nodo
precedente e al nodo successivo. A differenza del \inlinecode{vector},
gli elementi non sono contigui in memoria ma allocati separatamente
nello heap e collegati tra loro tramite i puntatori.

\subsection{Dichiarazione e dimensione}

\begin{itemize}
    \item \inlinecode{list<int> l;} --- crea una lista vuota di interi.
    \item \inlinecode{list<int> l(n);} --- crea una lista con \inlinecode{n}
          elementi inizializzati al valore di default (\inlinecode{0} per
          \inlinecode{int}).
    \item \inlinecode{l.size();} --- restituisce il numero di elementi
          presenti nella lista.
\end{itemize}

\subsection{Inserimento e rimozione in testa e in coda}

La \inlinecode{list} supporta inserimento e rimozione in tempo $O(1)$
sia in testa che in coda:

\begin{itemize}
    \item \inlinecode{l.push\_back(x);} --- aggiunge l'elemento \inlinecode{x}
          in fondo alla lista.
    \item \inlinecode{l.push\_front(x);} --- aggiunge l'elemento \inlinecode{x}
          in testa alla lista.
    \item \inlinecode{l.pop\_back();} --- rimuove l'ultimo elemento.
    \item \inlinecode{l.pop\_front();} --- rimuove il primo elemento.
\end{itemize}












\section{Problemi e pattern di programmazione }


\begin{problema}
\textit{Funzione di lettura di un vettore.}\\
Scrivere una funzione \inlinecode{leggi} che riceva un vettore di interi per riferimento,
legga da standard input il numero di elementi e li carichi uno ad uno nel vettore.
\end{problema}

\strategia{Leggere prima la dimensione $n$, ridimensionare il vettore con \inlinecode{v.resize(n)},
poi caricare i valori uno ad uno sfruttando il range-for con riferimento.}

\begin{soluzione}
Un pattern comune \`e definire una funzione \inlinecode{leggi} che riceve
un vettore per riferimento, legge da standard input il numero di elementi
e poi li carica uno ad uno nel vettore:
\begin{enumerate}
  \item Legge un intero $n$ dallo standard input, che rappresenta
  il numero di elementi del vettore.
  \item Ridimensiona il vettore con \inlinecode{v.resize(n)}.
  \item Legge gli $n$ valori uno ad uno, caricandoli direttamente
  negli elementi del vettore.
\end{enumerate}
\begin{codeexample}[Funzione \lstinline{leggi}]
  \begin{codebox}
    \begin{lstlisting}
void leggi(std::vector<int>& v) {
    int n;
    cin >> n;
    v.resize(n);
    for (int& x : v)
        cin >> x;
}
    \end{lstlisting}
  \end{codebox}
\end{codeexample}
Il ciclo range-for con \inlinecode{int\& x} dichiara \inlinecode{x}
come \textbf{reference} all'elemento corrente del vettore.
A ogni iterazione \inlinecode{x} si lega successivamente a ciascun
elemento di \inlinecode{v}: poich\`e \`e una reference,
l'istruzione \inlinecode{cin >> x} scrive direttamente
nell'elemento del vettore, modificandolo in place.
\end{soluzione}

\subsection{Pattern esistenziale: \'e presente?}

Un pattern ricorrente sui vettori \`e verificare se esiste almeno un elemento
che soddisfa una certa propriet\`a --- si parla di \textbf{propriet\`a esistenziale}.
L'idea \`e usare un flag booleano \inlinecode{found} inizializzato a \inlinecode{false},
che viene impostato a \inlinecode{true} non appena si trova un elemento che soddisfa la condizione;
il ciclo termina anticipatamente grazie alla condizione \inlinecode{\&\& !found}.

\begin{problema}
\textit{Presenza di un elemento.}\\
Data una funzione \inlinecode{presente} che riceva un vettore di interi
\inlinecode{const vector<int>\& v} e un intero \inlinecode{x},
decidere se \inlinecode{x} \`e presente nel vettore.
\end{problema}

\strategia{Usare un flag \inlinecode{found = false} e scorrere il vettore
fino a trovare l'elemento oppure fino alla fine; restituire il flag.}

\begin{soluzione}
\begin{codeexample}[Funzione \lstinline{presente}]
  \begin{codebox}
    \begin{lstlisting}
bool presente(const vector<int>& v, int x) {
    bool found = false;
    int i = 0;
    while (i < v.size() && !found) {
        if (v.at(i) == x)
            found = true;
        i++;
    }
    return found;
}
    \end{lstlisting}
  \end{codebox}
\end{codeexample}
\end{soluzione}










\pagebreak 

\subsection{Pattern universale: tutti soddisfano la propriet\`a?}

Il duale della propriet\`a esistenziale \`e la \textbf{propriet\`a universale}:
vogliamo verificare se \emph{tutti} gli elementi soddisfano una condizione.
In questo caso il flag parte a \inlinecode{true} e viene abbattuto
non appena si trova un elemento che \emph{non} soddisfa la condizione.

\begin{problema}
\textit{Tutti i pari.}\\
Data la funzione \inlinecode{tuttiPari} che riceva un vettore di interi,
decidere se tutti gli elementi sono pari.
\end{problema}

\strategia{Usare un flag \inlinecode{found = true} e scorrere il vettore;
se si incontra un elemento dispari abbattere il flag e interrompere il ciclo.}

\begin{soluzione}
\begin{codeexample}[Funzione \lstinline{tuttiPari}]
  \begin{codebox}
    \begin{lstlisting}
bool tuttiPari(const vector<int>& v) {
    bool found = true;
    int i = 0;
    while (i < v.size() && found) {
        if (v.at(i) % 2 != 0)
            found = false;
        i++;
    }
    return found;
}
    \end{lstlisting}
  \end{codebox}
\end{codeexample}
\end{soluzione}


















\begin{problema}
\textit{Implementazione del Bubble Sort.}\\
Scrivere una funzione che ordini un vettore di interi in ordine crescente usando l'algoritmo bubble sort.
\end{problema}

\strategia{Dividere il vettore in due zone: da 0 a \texttt{soglia - 1} ordinato, da \texttt{soglia} a fine disordinato. Ad ogni passata, l'elemento più piccolo della zona disordinata ``galleggia'' verso la posizione \texttt{soglia} attraverso scambi ripetuti, poi \texttt{soglia} avanza.}

\begin{soluzione}
Il bubble sort funziona facendo ``galleggiare'' gli elementi nella posizione corretta
attraverso scambi ripetuti tra elementi adiacenti.
Abbiamo un vettore diviso in due zone da una variabile \inlinecode{soglia}:
da indice 0 a \inlinecode{soglia - 1} la zona \`e ordinata, da \inlinecode{soglia}
a fine vettore la zona \`e disordinata.
All'inizio \inlinecode{soglia = 0} (niente \`e ordinato).
Ad ogni passata, l'elemento pi\`u piccolo della zona disordinata ``galleggia''
fino alla posizione \inlinecode{soglia}, e poi \inlinecode{soglia} avanza di 1.

\begin{codeexample}[Funzione \lstinline{bubble}]
  \begin{codebox}
    \begin{lstlisting}
void bubble(vector<int>& v) {
    if (v.empty()) return;          // (1)
    int soglia = 0;                 // (2)
    while (soglia < (int)v.size() - 1) {   // (3)
        for (int i = v.size() - 1; i > soglia; i--) {  // (4)
            if (v.at(i) < v.at(i - 1)) {         // (5)
                swap(v.at(i), v.at(i - 1));      // (6)
            }
        }
        soglia++;                   // (7)
    }
}
    \end{lstlisting}
  \end{codebox}
\end{codeexample}

\textbf{Analisi riga per riga:}
\begin{enumerate}
  \item \inlinecode{if (v.empty()) return;} --- Se il vettore \`e vuoto non c'\`e nulla da ordinare.
  Senza questo controllo, \inlinecode{v.size() - 1} su un vettore vuoto darebbe un numero enorme
  (perch\'e \inlinecode{size()} restituisce un \inlinecode{unsigned}, e $0 - 1$ in unsigned va in overflow).
  
  \item \inlinecode{int soglia = 0;} --- Inizializziamo la soglia a 0:
  ``nessun elemento \`e ancora nella posizione corretta''.
  Tutta la zona da 0 a \inlinecode{v.size()-1} \`e disordinata.
  
  \item \inlinecode{while (soglia < (int)v.size() - 1)} --- Il ciclo esterno continua finch\'e
  la soglia non ha raggiunto la penultima posizione.
  Perch\'e penultima e non ultima? Perch\'e se tutti gli elementi tranne l'ultimo sono al posto giusto,
  anche l'ultimo per forza lo \`e.
  Il cast \inlinecode{(int)} serve perch\'e \inlinecode{v.size()} \`e unsigned.
  
  \item \inlinecode{for (int i = v.size() - 1; i > soglia; i--)} --- Questo \`e il ciclo interno.
  Parte dall'ultimo elemento del vettore e scende verso la soglia.
  Perch\'e dal fondo? Perch\'e vogliamo far ``salire'' l'elemento pi\`u piccolo verso sinistra,
  verso la posizione soglia.
  \inlinecode{i > soglia} e non \inlinecode{i >= soglia} perch\'e nel corpo del \inlinecode{for}
  confrontiamo \inlinecode{v[i]} con \inlinecode{v[i-1]}.
  Se \inlinecode{i} fosse uguale a \inlinecode{soglia}, allora \inlinecode{i-1} sarebbe
  \inlinecode{soglia - 1}, cio\`e andremmo a toccare la zona gi\`a ordinata, che non va toccata.
  
  \item \inlinecode{if (v.at(i) < v.at(i - 1))} --- Confrontiamo l'elemento corrente con quello che lo precede.
  Se quello corrente \`e pi\`u piccolo, sono nell'ordine sbagliato (perch\'e vogliamo ordine crescente).
  
  \item \inlinecode{swap(v.at(i), v.at(i - 1));} --- Li scambiamo.
  Dopo lo scambio, l'elemento pi\`u piccolo si \`e spostato di una posizione verso sinistra.
  Questo succede ripetutamente lungo tutto il ciclo \inlinecode{for},
  quindi l'elemento pi\`u piccolo della zona disordinata ``galleggia'' fino a \inlinecode{soglia}.
  
  \item \inlinecode{soglia++;} --- La passata \`e finita.
  L'elemento pi\`u piccolo della zona disordinata \`e arrivato in posizione \inlinecode{soglia}.
  Ora quella posizione fa parte della zona ordinata, quindi avanziamo la soglia di 1.
\end{enumerate}
\end{soluzione}

\subsection{Esempio di esecuzione}

\section*{Legenda}
\begin{tikzpicture}
    \node[draw, fill=ordinato, minimum width=0.6cm, minimum height=0.6cm, rounded corners=2pt] at (0,0) {};
    \node[right] at (0.5,0) {Ordinato (sotto soglia)};
    \node[draw, fill=confronto, minimum width=0.6cm, minimum height=0.6cm, rounded corners=2pt] at (5.4,0) {};
    \node[right] at (5.8,0) {Confronto};
    \node[draw, fill=scambio, minimum width=0.6cm, minimum height=0.6cm, rounded corners=2pt] at (9,0) {};
    \node[right] at (9.5,0) {Scambio};
    \node[draw, fill=neutro, minimum width=0.6cm, minimum height=0.6cm, rounded corners=2pt] at (13,0) {};
    \node[right] at (13.5,0) {Disordinato};
\end{tikzpicture}

\bigskip
Vettore iniziale: $[5, 3, 1, 4, 2]$, \quad $\text{soglia} = 0$

%===========================================
\section*{Passata 1}
$\text{soglia} = 0$, \quad il \texttt{for} va da $i=4$ fino a $i>0$ \quad (4 confronti)

\bigskip
\textbf{Stato iniziale:}
\begin{center}
\begin{tikzpicture}
    \foreach \x/\v in {0/5, 1.2/3, 2.4/1, 3.6/4, 4.8/2} {
        \cella{\x}{neutro}; \num{\x}{\v};
    }
\end{tikzpicture}
\end{center}

$i=4$: $v[4]=2 < v[3]=4$ \checkmark $\Rightarrow$ \textbf{swap!}
\begin{center}
\begin{tikzpicture}
    \foreach \x/\v/\c in {0/5/neutro, 1.2/3/neutro, 2.4/1/neutro, 3.6/2/scambio, 4.8/4/scambio} {
        \cella{\x}{\c}; \num{\x}{\v};
    }
\end{tikzpicture}
\end{center}

$i=3$: $v[3]=2 < v[2]=1$? No $\Rightarrow$ nessuno scambio
\begin{center}
\begin{tikzpicture}
    \foreach \x/\v/\c in {0/5/neutro, 1.2/3/neutro, 2.4/1/confronto, 3.6/2/confronto, 4.8/4/neutro} {
        \cella{\x}{\c}; \num{\x}{\v};
    }
\end{tikzpicture}
\end{center}

$i=2$: $v[2]=1 < v[1]=3$ \checkmark $\Rightarrow$ \textbf{swap!}
\begin{center}
\begin{tikzpicture}
    \foreach \x/\v/\c in {0/5/neutro, 1.2/1/scambio, 2.4/3/scambio, 3.6/2/neutro, 4.8/4/neutro} {
        \cella{\x}{\c}; \num{\x}{\v};
    }
\end{tikzpicture}
\end{center}

$i=1$: $v[1]=1 < v[0]=5$ \checkmark $\Rightarrow$ \textbf{swap!}
\begin{center}
\begin{tikzpicture}
    \foreach \x/\v/\c in {0/1/scambio, 1.2/5/scambio, 2.4/3/neutro, 3.6/2/neutro, 4.8/4/neutro} {
        \cella{\x}{\c}; \num{\x}{\v};
    }
\end{tikzpicture}
\end{center}

$\text{soglia++} \Rightarrow \text{soglia} = 1$
\begin{center}
\begin{tikzpicture}
    \foreach \x/\v/\c in {0/1/ordinato, 1.2/5/neutro, 2.4/3/neutro, 3.6/2/neutro, 4.8/4/neutro} {
        \cella{\x}{\c}; \num{\x}{\v};
    }
\end{tikzpicture}
\end{center}

%===========================================
\section*{Passata 2}
$\text{soglia} = 1$, \quad il \texttt{for} va da $i=4$ fino a $i>1$ \quad (3 confronti)

\bigskip
$i=4$: $v[4]=4 < v[3]=2$? No $\Rightarrow$ nessuno scambio
\begin{center}
\begin{tikzpicture}
    \foreach \x/\v/\c in {0/1/ordinato, 1.2/5/neutro, 2.4/3/neutro, 3.6/2/confronto, 4.8/4/confronto} {
        \cella{\x}{\c}; \num{\x}{\v};
    }
\end{tikzpicture}
\end{center}

$i=3$: $v[3]=2 < v[2]=3$ \checkmark $\Rightarrow$ \textbf{swap!}
\begin{center}
\begin{tikzpicture}
    \foreach \x/\v/\c in {0/1/ordinato, 1.2/5/neutro, 2.4/2/scambio, 3.6/3/scambio, 4.8/4/neutro} {
        \cella{\x}{\c}; \num{\x}{\v};
    }
\end{tikzpicture}
\end{center}

$i=2$: $v[2]=2 < v[1]=5$ \checkmark $\Rightarrow$ \textbf{swap!}
\begin{center}
\begin{tikzpicture}
    \foreach \x/\v/\c in {0/1/ordinato, 1.2/2/scambio, 2.4/5/scambio, 3.6/3/neutro, 4.8/4/neutro} {
        \cella{\x}{\c}; \num{\x}{\v};
    }
\end{tikzpicture}
\end{center}

$\text{soglia++} \Rightarrow \text{soglia} = 2$
\begin{center}
\begin{tikzpicture}
    \foreach \x/\v/\c in {0/1/ordinato, 1.2/2/ordinato, 2.4/5/neutro, 3.6/3/neutro, 4.8/4/neutro} {
        \cella{\x}{\c}; \num{\x}{\v};
    }
\end{tikzpicture}
\end{center}

%===========================================
\section*{Passata 3}
$\text{soglia} = 2$, \quad il \texttt{for} va da $i=4$ fino a $i>2$ \quad (2 confronti)

\bigskip
$i=4$: $v[4]=4 < v[3]=3$? No $\Rightarrow$ nessuno scambio
\begin{center}
\begin{tikzpicture}
    \foreach \x/\v/\c in {0/1/ordinato, 1.2/2/ordinato, 2.4/5/neutro, 3.6/3/confronto, 4.8/4/confronto} {
        \cella{\x}{\c}; \num{\x}{\v};
    }
\end{tikzpicture}
\end{center}

$i=3$: $v[3]=3 < v[2]=5$ \checkmark $\Rightarrow$ \textbf{swap!}
\begin{center}
\begin{tikzpicture}
    \foreach \x/\v/\c in {0/1/ordinato, 1.2/2/ordinato, 2.4/3/scambio, 3.6/5/scambio, 4.8/4/neutro} {
        \cella{\x}{\c}; \num{\x}{\v};
    }
\end{tikzpicture}
\end{center}

$\text{soglia++} \Rightarrow \text{soglia} = 3$
\begin{center}
\begin{tikzpicture}
    \foreach \x/\v/\c in {0/1/ordinato, 1.2/2/ordinato, 2.4/3/ordinato, 3.6/5/neutro, 4.8/4/neutro} {
        \cella{\x}{\c}; \num{\x}{\v};
    }
\end{tikzpicture}
\end{center}

%===========================================
\section*{Passata 4}
$\text{soglia} = 3$, \quad il \texttt{for} va da $i=4$ fino a $i>3$ \quad (1 confronto)

\bigskip
$i=4$: $v[4]=4 < v[3]=5$ \checkmark $\Rightarrow$ \textbf{swap!}
\begin{center}
\begin{tikzpicture}
    \foreach \x/\v/\c in {0/1/ordinato, 1.2/2/ordinato, 2.4/3/ordinato, 3.6/4/scambio, 4.8/5/scambio} {
        \cella{\x}{\c}; \num{\x}{\v};
    }
\end{tikzpicture}
\end{center}

$\text{soglia++} \Rightarrow \text{soglia} = 4$
\begin{center}
\begin{tikzpicture}
    \foreach \x/\v/\c in {0/1/ordinato, 1.2/2/ordinato, 2.4/3/ordinato, 3.6/4/ordinato, 4.8/5/ordinato} {
        \cella{\x}{\c}; \num{\x}{\v};
    }
\end{tikzpicture}
\end{center}

%===========================================
\section*{Terminazione}

$\text{soglia} = 4 < v.\text{size}() - 1 = 4$ ? \quad \textbf{No} $\Rightarrow$ il \texttt{while} termina.

\bigskip
\textbf{Risultato finale:} $[1, 2, 3, 4, 5]$ \checkmark



\subsection{Ottimizzazione: terminazione anticipata}

\begin{problema}
\textit{Bubble Sort con early termination.}\\
Modificare la funzione bubble sort per terminare anticipatamente se il vettore è già ordinato.
\end{problema}

\strategia{Usare un flag booleano \texttt{check} che viene resettato a \texttt{false} all'inizio di ogni passata. Se durante la passata avviene almeno uno scambio, \texttt{check} viene impostato a \texttt{true}. Se alla fine di una passata \texttt{check} è ancora \texttt{false}, significa che non ci sono stati scambi e il vettore è già ordinato, quindi possiamo terminare.}

\begin{soluzione}
\begin{codeexample}[Bubble sort con early termination]
  \begin{codebox}
    \begin{lstlisting}
void bubble(vector<int>& v) {
    if (v.empty()) return;

    bool check = true;
    int soglia = 0;

    while (soglia < (int)v.size() - 1 && check) {
        check = false;
        for (int i = v.size() - 1; i > soglia; i--) {
            if (v[i] < v[i - 1]) {
                swap(v[i], v[i - 1]);
                check = true;
            }
        }
        soglia++;
    }
}
    \end{lstlisting}
  \end{codebox}
\end{codeexample}
\end{soluzione}

\subsection{Esempio di early termination}

\section*{Legenda}
\begin{tikzpicture}
    \node[draw, fill=confronto, minimum width=0.6cm, minimum height=0.6cm, rounded corners=2pt] at (0,0) {};
    \node[right] at (0.5,0) {Confronto};
    \node[draw, fill=scambio, minimum width=0.6cm, minimum height=0.6cm, rounded corners=2pt] at (4,0) {};
    \node[right] at (4.5,0) {Scambio};
    \node[draw, fill=neutro, minimum width=0.6cm, minimum height=0.6cm, rounded corners=2pt] at (8,0) {};
    \node[right] at (8.5,0) {Non toccato};
    \node[draw, fill=ordinato, minimum width=0.6cm, minimum height=0.6cm, rounded corners=2pt] at (12,0) {};
    \node[right] at (12.5,0) {Tutto ordinato};
\end{tikzpicture}

\bigskip
\textbf{Nota:} in questo esercizio \textbf{non c'\`e la soglia}. Il \texttt{for} scorre sempre tutto il vettore. L'unica ottimizzazione \`e il flag \texttt{check}: se una passata intera non produce scambi, il vettore \`e gi\`a ordinato e ci fermiamo.

\bigskip
Vettore iniziale: $[1, 3, 2, 4, 5]$ (quasi ordinato, solo 3 e 2 invertiti)

%===========================================
\section*{Passata 1}
\texttt{check} viene resettato a \texttt{false} \quad $\Rightarrow$ \quad il \texttt{for} va da $i=4$ fino a $i>0$ \quad (4 confronti)

\bigskip
\textbf{Stato iniziale:}
\begin{center}
\begin{tikzpicture}
    \foreach \x/\v in {0/1, 1.2/3, 2.4/2, 3.6/4, 4.8/5} {
        \cella{\x}{neutro}; \num{\x}{\v};
    }
\end{tikzpicture}
\end{center}

$i=4$: $v[4]=5 < v[3]=4$? No $\Rightarrow$ nessuno scambio
\begin{center}
\begin{tikzpicture}
    \foreach \x/\v/\c in {0/1/neutro, 1.2/3/neutro, 2.4/2/neutro, 3.6/4/confronto, 4.8/5/confronto} {
        \cella{\x}{\c}; \num{\x}{\v};
    }
\end{tikzpicture}
\end{center}

$i=3$: $v[3]=4 < v[2]=2$? No $\Rightarrow$ nessuno scambio
\begin{center}
\begin{tikzpicture}
    \foreach \x/\v/\c in {0/1/neutro, 1.2/3/neutro, 2.4/2/confronto, 3.6/4/confronto, 4.8/5/neutro} {
        \cella{\x}{\c}; \num{\x}{\v};
    }
\end{tikzpicture}
\end{center}

$i=2$: $v[2]=2 < v[1]=3$ \checkmark $\Rightarrow$ \textbf{swap!} \quad \texttt{check = true}
\begin{center}
\begin{tikzpicture}
    \foreach \x/\v/\c in {0/1/neutro, 1.2/2/scambio, 2.4/3/scambio, 3.6/4/neutro, 4.8/5/neutro} {
        \cella{\x}{\c}; \num{\x}{\v};
    }
\end{tikzpicture}
\end{center}

$i=1$: $v[1]=2 < v[0]=1$? No $\Rightarrow$ nessuno scambio
\begin{center}
\begin{tikzpicture}
    \foreach \x/\v/\c in {0/1/confronto, 1.2/2/confronto, 2.4/3/neutro, 3.6/4/neutro, 4.8/5/neutro} {
        \cella{\x}{\c}; \num{\x}{\v};
    }
\end{tikzpicture}
\end{center}

\textbf{Fine passata 1:} \texttt{check = true} $\Rightarrow$ \`e avvenuto almeno uno scambio $\Rightarrow$ \textbf{si continua}

\begin{center}
\begin{tikzpicture}
    \foreach \x/\v in {0/1, 1.2/2, 2.4/3, 3.6/4, 4.8/5} {
        \cella{\x}{neutro}; \num{\x}{\v};
    }
\end{tikzpicture}
\end{center}

%===========================================
\section*{Passata 2}
\texttt{check} viene resettato a \texttt{false} \quad $\Rightarrow$ \quad il \texttt{for} va da $i=4$ fino a $i>0$ \quad (4 confronti)

\bigskip
$i=4$: $v[4]=5 < v[3]=4$? No $\Rightarrow$ nessuno scambio
\begin{center}
\begin{tikzpicture}
    \foreach \x/\v/\c in {0/1/neutro, 1.2/2/neutro, 2.4/3/neutro, 3.6/4/confronto, 4.8/5/confronto} {
        \cella{\x}{\c}; \num{\x}{\v};
    }
\end{tikzpicture}
\end{center}

$i=3$: $v[3]=4 < v[2]=3$? No $\Rightarrow$ nessuno scambio
\begin{center}
\begin{tikzpicture}
    \foreach \x/\v/\c in {0/1/neutro, 1.2/2/neutro, 2.4/3/confronto, 3.6/4/confronto, 4.8/5/neutro} {
        \cella{\x}{\c}; \num{\x}{\v};
    }
\end{tikzpicture}
\end{center}

$i=2$: $v[2]=3 < v[1]=2$? No $\Rightarrow$ nessuno scambio
\begin{center}
\begin{tikzpicture}
    \foreach \x/\v/\c in {0/1/neutro, 1.2/2/confronto, 2.4/3/confronto, 3.6/4/neutro, 4.8/5/neutro} {
        \cella{\x}{\c}; \num{\x}{\v};
    }
\end{tikzpicture}
\end{center}

$i=1$: $v[1]=2 < v[0]=1$? No $\Rightarrow$ nessuno scambio
\begin{center}
\begin{tikzpicture}
    \foreach \x/\v/\c in {0/1/confronto, 1.2/2/confronto, 2.4/3/neutro, 3.6/4/neutro, 4.8/5/neutro} {
        \cella{\x}{\c}; \num{\x}{\v};
    }
\end{tikzpicture}
\end{center}

\textbf{Fine passata 2:} \texttt{check = false} $\Rightarrow$ \textbf{nessuno scambio in tutta la passata!}

\bigskip
Questo significa che per ogni $i$: $v[i-1] \leq v[i]$, cio\`e ogni elemento \`e minore o uguale del successivo.

Quindi il vettore \`e \textbf{gi\`a ordinato} e possiamo \textbf{terminare}.

\begin{center}
\begin{tikzpicture}
    \foreach \x/\v in {0/1, 1.2/2, 2.4/3, 3.6/4, 4.8/5} {
        \cella{\x}{ordinato}; \num{\x}{\v};
    }
\end{tikzpicture}
\end{center}

%===========================================
\section*{Terminazione}

Il \texttt{while} controlla \texttt{check}:

$$\texttt{check} = \texttt{false} \quad \Rightarrow \quad \text{il \texttt{while} termina}$$

\textbf{Risultato finale:} $[1, 2, 3, 4, 5]$ \checkmark

%===========================================
\section*{Confronto con il Bubble Sort senza early termination}

\begin{center}
\begin{tabular}{|l|c|c|}
\hline
& \textbf{Senza early termination} & \textbf{Con early termination} \\
\hline
Passate eseguite & $n - 1 = 4$ & $2$ \\
\hline
Confronti totali & $\dfrac{n(n-1)}{2} = 10$ & $4 + 4 = 8$ \\
\hline
\end{tabular}
\end{center}

\bigskip
In questo esempio abbiamo risparmiato 2 passate e 2 confronti. Il vantaggio diventa enorme con vettori grandi e quasi ordinati.

%===========================================
\section*{Caso estremo: vettore gi\`a ordinato}

Vettore: $[1, 2, 3, 4, 5]$

\bigskip
\textbf{Passata 1:} \texttt{check} resettato a \texttt{false}

\begin{center}
\begin{tabular}{ll}
$i=4$: $v[4]=5 < v[3]=4$? & No \\
$i=3$: $v[3]=4 < v[2]=3$? & No \\
$i=2$: $v[2]=3 < v[1]=2$? & No \\
$i=1$: $v[1]=2 < v[0]=1$? & No \\
\end{tabular}
\end{center}

\texttt{check} \`e ancora \texttt{false} $\Rightarrow$ \textbf{esco dopo una sola passata!}

$$\text{Costo:} \quad n - 1 = 4 \text{ confronti} \quad \Rightarrow \quad O(n)$$

Senza early termination avremmo fatto $\dfrac{n(n-1)}{2} = 10$ confronti $\Rightarrow O(n^2)$.

%===========================================
\section*{Riepilogo complessit\`a}

\begin{align*}
    \textbf{Caso migliore} \text{ (vettore gi\`a ordinato):} &\quad O(n) \\
    \textbf{Caso medio:} &\quad O(n^2) \\
    \textbf{Caso peggiore} \text{ (vettore ordinato al contrario):} &\quad O(n^2) \\
    \textbf{Spazio:} &\quad O(1) \text{ (in-place)}
\end{align*}



\chapter{Ricorsione}
\intro{}


\section{Introduzione}

\defn{Ricorsione}{
    La ricorsione è un metodo di risoluzione di problemi computazionali in cui la soluzione di un problema dipende dalla soluzione di istanze più piccole dello stesso problema 
    e questo processo si realizza attraverso una funzione che richiama se stessa dall'interno del proprio corpo.
}

Una funzione ricorsiva è corretta e completa è sempre composta da esattamente due componenti fondamentali e inscindibili:

\begin{itemize}
    \item Caso base: il caso più elementare del problema, per cui la risposta è nota e immediata  la funzione lo risolve senza richiamare se stessa. È la condizione di arresto che garantisce la terminazione della catena ricorsiva. Senza di esso, la ricorsione continuerebbe all'infinito.
    \item Passo ricorsivo: il caso generale, in cui il problema non è ancora nella sua forma più semplice La funzione lo riduce in una versione più piccola e si richiama su quella versione ridotta, avanzando progressivamente verso il caso base.
\end{itemize}
\subsection{Struttura formale}
Matematicamente, una definizione ricorsiva si esprime sempre in questa forma:
\[
\begin{cases}
    \text{soluzione diretta} \ \text{se } n=\text{caso base}\\
    g\bigl(n,f(n')\bigr) \qquad \quad \text{altrimenti , dove } n'<n
\end{cases}
\]

Dove $n'$ è una versione ridotta di $n$, e $g$ è una qualsiasi operazione che combina il risultato della chiamata ricorsiva con il valore corrente. L'importante è che ogni passo ricorsivo garantisca che $n'$ converga verso il  caso base.


Da ora consideriamo la ricorsione non come “trucco”, ma come un insieme di pattern mentali ricorrenti.

\section{Ricorsione lineare}
\defn{Ricorsone lineare}{
    Una ricorsione è \textbf{lineare} quando ogni chiamataata genera esattamente una chiamata ricorsiva. La struttura è una catena: il problema di taglia n viene ridotto a un problema della stessa natura ma più piccolo, fino al caso base
}
Il principio è identico all'induzione matematica: il caso base corrisponde al passo base dell'induzione il passo ricorsivo corrisponde al passo induttivo

Non va simulata la catena nella testa.  Si Assume che la funzione funzioni già per input più piccoli.

\newpage

La ricorsione lineare si divide in tre sottotipi a seconda di dove avviene il lavoro rispetto alla chiamata ricorsiva.

\subsection{Post-order}
    Il lavoro avviene \textbf{dopo} la chiamata ricorsiva. La chiamata ricorsiva viene fatta prima, il suo risultato viene usato dopo.


\begin{codeexample}[Somma elementi positivi]
  \strategia{
    \begin{itemize}
      \item Facciamo la \textbf{ricorsione sull'indice} \inlinecode{i},
            avanzato di 1 ad ogni passo.
      \item \textbf{Ipotesi induttiva:} la funzione è corretta per qualsiasi
            indice strettamente maggiore di \inlinecode{i}.
      \item \textbf{Passo ricorsivo:} uso la funzione (corretta per ipotesi
            su \inlinecode{i+1}) e aggiungo il contributo di \inlinecode{v.at(i)}.
            \vspace{0.1cm}

            \inlinecode{int contributo = v.at(i) > 0 ? v.at(i) : 0;}
            \vspace{0.1cm}

            \inlinecode{return contributo + somma(v, i + 1);}

      \item \textbf{Caso base:} quando \inlinecode{i == (int)v.size()}
            non ci sono più elementi da esaminare; la somma è 0.
            \vspace{0.1cm}

            \inlinecode{if (i == (int)v.size()) return 0;}
    \end{itemize}
  }

  \begin{codebox}
    \begin{lstlisting}
int somma(const vector<int>& v, int i = 0) {
    if (i == (int)v.size()) return 0;
    int contributo = v.at(i) > 0 ? v.at(i) : 0;
    return contributo + somma(v, i + 1);
}
    \end{lstlisting}
  \end{codebox}

\end{codeexample}



\begin{codeexample}[Massimo]
  \strategia{
    \begin{itemize}
      \item Facciamo la \textbf{ricorsione sull'indice} \inlinecode{i},
            avanzato di 1 ad ogni passo.
      \item \textbf{Ipotesi induttiva:} la funzione è corretta per qualsiasi
            indice strettamente maggiore di \inlinecode{i}.
      \item \textbf{Passo ricorsivo:} uso la funzione (corretta per ipotesi
            su \inlinecode{i+1}) e confronto \inlinecode{v.at(i)} con il
            massimo del resto già noto.
            \vspace{0.1cm}

            \inlinecode{return max(v.at(i), massimo(v, i + 1));}

      \item \textbf{Caso base:} quando \inlinecode{i == (int)v.size() - 1}
            c'è un solo elemento, il massimo è lui.
            \vspace{0.1cm}

            \inlinecode{if (i == (int)v.size() - 1) return v.at(i);}
    \end{itemize}
  }

  \begin{codebox}
    \begin{lstlisting}
int massimo(const vector<int>& v, int i = 0) {
    if (i == (int)v.size() - 1) return v.at(i);
    return max(v.at(i), massimo(v, i + 1));
}
    \end{lstlisting}
  \end{codebox}

\end{codeexample}

  \subsection{Pre-order}
Il lavoro avviene prima della chiamata ricorsiva. L'effetto è l'opposto della post-order: si elabora in ordine naturale invece che inverso.
  \begin{codeexample}[Stampa da 1 a n]
  \strategia{
    \begin{itemize}
      \item Facciamo la \textbf{ricorsione sull'indice} \inlinecode{i},
            avanzato di 1 ad ogni passo.
      \item \textbf{Ipotesi induttiva:} la funzione è corretta per qualsiasi
            indice strettamente maggiore di \inlinecode{i}.
      \item \textbf{Passo ricorsivo:} uso la funzione (corretta per ipotesi
            su \inlinecode{i+1}) e stampo \inlinecode{i} prima di delegare
            il resto, il lavoro avviene alla discesa (pre-order).
            \vspace{0.1cm}

            \inlinecode{cout << i << " ";}
            \vspace{0.1cm}

            \inlinecode{stampa(i + 1, n);}

      \item \textbf{Caso base:} quando \inlinecode{i > n} non c'è più
            nulla da stampare.
            \vspace{0.1cm}

            \inlinecode{if (i > n) return;}
    \end{itemize}
  }

  \begin{codebox}
    \begin{lstlisting}
void stampa(int i, int n) {
    if (i > n) return;
    cout << i << " ";
    stampa(i + 1, n);
}
    \end{lstlisting}
  \end{codebox}

\end{codeexample}




















\begin{codeexample}[Somma con accumulatore]
  \strategia{
    \begin{itemize}
      \item Facciamo la \textbf{ricorsione sull'indice} \inlinecode{i},
            avanzato di 1 ad ogni passo, con \inlinecode{acc} che porta
            la somma parziale degli elementi positivi già visitati.
      \item \textbf{Ipotesi induttiva:} la funzione è corretta per qualsiasi
            indice strettamente maggiore di \inlinecode{i}.
      \item \textbf{Passo ricorsivo:} uso la funzione (corretta per ipotesi
            su \inlinecode{i+1}) aggiornando l'accumulatore con il contributo
            dell'elemento corrente  nessun calcolo rimane pendente al ritorno.
            \vspace{0.1cm}

            \inlinecode{int contributo = v.at(i) > 0 ? v.at(i) : 0;}
            \vspace{0.1cm}

            \inlinecode{return somma(v, i + 1, acc + contributo);}

      \item \textbf{Caso base:} quando \inlinecode{i == v.size()} tutti gli
            elementi sono stati processati; viene restituito direttamente
            l'accumulatore \inlinecode{acc}.
            \vspace{0.1cm}

            \inlinecode{if (i == v.size()) return acc;}
    \end{itemize}
  }

  \begin{codebox}
    \begin{lstlisting}
int somma(const vector<int>& v, int i = 0, int acc = 0) {
    if (i == v.size()) return acc;
    int contributo = v.at(i) > 0 ? v.at(i) : 0;
    return somma(v, i + 1, acc + contributo);
}
    \end{lstlisting}
  \end{codebox}

\end{codeexample}
  \subsection{Tail recursion}
  La chiamata ricorsiva è l'ultima operazione assoluta della funzione nessun calcolo pendente dopo di essa. Lo stato che nella post-order si accumula nello stack viene qui portato in avanti come parametro accumulatore.

  \begin{codeexample}[Somma con accumulatore]
  \strategia{
    \begin{itemize}
      \item Facciamo la \textbf{ricorsione sull'indice} \inlinecode{i},
            avanzato di 1 ad ogni passo, con \inlinecode{acc} che porta
            la somma parziale degli elementi positivi già visitati.
      \item \textbf{Ipotesi induttiva:} la funzione è corretta per qualsiasi
            indice strettamente maggiore di \inlinecode{i}.
      \item \textbf{Passo ricorsivo:} uso la funzione (corretta per ipotesi
            su \inlinecode{i+1}) aggiornando l'accumulatore con il contributo
            dell'elemento corrente  nessun calcolo rimane pendente al ritorno.
            \vspace{0.1cm}

            \inlinecode{int contributo = v.at(i) > 0 ? v.at(i) : 0;}
            \vspace{0.1cm}

            \inlinecode{return somma(v, i + 1, acc + contributo);}

      \item \textbf{Caso base:} quando \inlinecode{i == (int)v.size()} tutti
            gli elementi sono stati processati; viene restituito direttamente
            l'accumulatore \inlinecode{acc}.
            \vspace{0.1cm}

            \inlinecode{if (i == (int)v.size()) return acc;}
    \end{itemize}
  }

  \begin{codebox}
    \begin{lstlisting}
int somma(const vector<int>& v, int i = 0, int acc = 0) {
    if (i == (int)v.size()) return acc;
    int contributo = v.at(i) > 0 ? v.at(i) : 0;
    return somma(v, i + 1, acc + contributo);
}
    \end{lstlisting}
  \end{codebox}

\end{codeexample}

\begin{codeexample}[Massimo comun divisore di Euclide]
  La proprietà matematica su cui si fonda l'algoritmo è la seguente:
  \[
    \text{MCD}(a,\ b) = \text{MCD}(b,\ a \bmod b)
  \]
  \textbf{Perché è vera.} Dalla divisione euclidea si ha $a = q \cdot b + r$
  con $r = a \bmod b$. Qualsiasi divisore comune $d$ di $a$ e $b$ divide anche
  $r = a - q \cdot b$, e viceversa. I due insiemi di divisori comuni di
  $(a,b)$ e $(b,r)$ sono quindi identici, e in particolare lo è il loro massimo.

  \strategia{
    \begin{itemize}
      \item \textbf{Ricorsione su:} $b$, ridotto via $a \bmod b$ ad ogni passo.
            Poiché $0 \le a \bmod b < b$, il secondo argomento decresce
            strettamente la sequenza raggiunge $b = 0$ in tempo finito.
      \item \textbf{Ipotesi induttiva:} la funzione è corretta per qualsiasi
            valore di $b$ strettamente minore di quello corrente.
      \item \textbf{Passo ricorsivo:} per la proprietà sopra,
            $\text{MCD}(a,b) = \text{MCD}(b,\ a \bmod b)$;
            per ipotesi la funzione è già corretta su $(b,\ a \bmod b)$,
            quindi la chiamata ricorsiva restituisce il risultato corretto.
            La chiamata è \textit{tail}: $a \bmod b$ viene calcolato prima
            di essere passato, nessun calcolo rimane pendente al ritorno.
      \item \textbf{Caso base:} quando $b = 0$, ogni intero divide $0$,
            quindi $\text{MCD}(a,\ 0) = a$.
    \end{itemize}
  }

  \begin{codebox}
    \begin{lstlisting}
int mcd(int a, int b) {
    if (b == 0) return a;    
    return mcd(b, a % b);   
}
    \end{lstlisting}
  \end{codebox}

\end{codeexample}



Ora andiamo a vedere il secondo pattern la \textbf{Tail recursion}.

\section{Tail recursion}
\defn{Tail recursion}{
    Una chiamata ricorsiva è in coda, quando è l'ultima operazione assoluta della funzione nessun calcolo rimane pendente dopo di essa. Il valore restituito della chiama ricorsiva viene restituito direttamente ,senza ulteriori elaborazioni.
}

La tail recursion si divide in due famiglie principali: con accumulatore e senza accumulatore:

\subsection{Senza accumulatore}
La chiamata ricorsiva è  l'ultima operazione, ma non c'è nulla da comunicare in avanti ,  il risultato della chiamata ricorsiva viene restituito direttamente senza modifiche.

Si usa quando la ricorsione \textbf{elimina} possibilità invece di costruire qualcosa. Tipicamente ricerche, verifiche, riduzioni immediate. 


\begin{codeexample}[Ricerca lineare]
  \strategia{
    \begin{itemize}
      \item Facciamo la \textbf{ricorsione sull'indice} \inlinecode{i},
            avanzato di 1 ad ogni passo.
      \item \textbf{Ipotesi induttiva:} la funzione è corretta per qualsiasi
            indice strettamente maggiore di \inlinecode{i}.
      \item \textbf{Passo ricorsivo:} se \inlinecode{v.at(i)} non è il
            target, uso la funzione (corretta per ipotesi su \inlinecode{i+1})
            e ne restituisco direttamente il risultato  nessun calcolo
            rimane pendente al ritorno.
            \vspace{0.1cm}

            \inlinecode{return contiene(v, target, i + 1);}

      \item \textbf{Casi base:} due condizioni distinte.
            \vspace{0.1cm}

            \inlinecode{if (i == (int)v.size()) return false;}
            \vspace{0.1cm}

            \inlinecode{if (v.at(i) == target) return true;}
    \end{itemize}
  }

  \begin{codebox}
    \begin{lstlisting}
bool contiene(const vector<int>& v, int target, int i = 0) {
    if (i == (int)v.size()) return false;
    if (v.at(i) == target)  return true;
    return contiene(v, target, i + 1);
}
    \end{lstlisting}
  \end{codebox}

\end{codeexample}









\begin{codeexample}[Vettore ordinato]
  \strategia{
    \begin{itemize}
      \item Facciamo la \textbf{ricorsione sull'indice} \inlinecode{i},
            avanzato di 1 ad ogni passo.
      \item \textbf{Ipotesi induttiva:} la funzione è corretta per qualsiasi
            indice strettamente maggiore di \inlinecode{i}.
      \item \textbf{Passo ricorsivo:} se \inlinecode{v.at(i) <= v.at(i+1)},
            il confronto corrente è soddisfatto; uso la funzione (corretta
            per ipotesi su \inlinecode{i+1}) e ne restituisco direttamente
            il risultato  nessun calcolo rimane pendente al ritorno.
            \vspace{0.1cm}

            \inlinecode{return ordinato(v, i + 1);}

      \item \textbf{Casi base:} due condizioni distinte.
            \vspace{0.1cm}

            \inlinecode{if (v.at(i) > v.at(i + 1)) return false;}
            \vspace{0.1cm}

            \inlinecode{if (i == (int)v.size() - 1) return true;}
    \end{itemize}
  }

  \begin{codebox}
    \begin{lstlisting}
bool ordinato(const vector<int>& v, int i = 0) {
    if (i == (int)v.size() - 1)  return true; 
    if (v.at(i) > v.at(i + 1))   return false;
    return ordinato(v, i + 1);
}
    \end{lstlisting}
  \end{codebox}

\end{codeexample}


\subsection{Con accumulatore}
Serve quando si deve costruire o calcolare qualcosa lungo il percorso.
Questo sottotipo a sua volta ha altri sottotipi:

\subsubsection{Accumulatore singolo numerico}
Lo stato è un singolo valore numerico che si aggiorna ad ogni passo.




\begin{codeexample}[Fattoriale con accumulatore]

  Il \textbf{fattoriale} di un intero non-negativo $n$ ammette una
  definizione ricorsiva diretta:
  \[
    n! = \begin{cases}
      1              & \text{se } n = 0 \\
      n \cdot (n-1)! & \text{se } n > 0
    \end{cases}
  \]
  Per ottenere una \textbf{tail recursion}, si introduce un accumulatore
  \inlinecode{acc} che porta il prodotto parziale in avanti invece di
  lasciarlo pendente sullo stack. Poiché \inlinecode{acc} è un dettaglio
  implementativo interno, si separa la logica in due funzioni: una
  \textbf{funzione helper} che espone \inlinecode{acc} come parametro
  esplicito, e una \textbf{funzione wrapper} pubblica che la chiama
  inizializzando \inlinecode{acc} al valore neutro della moltiplicazione,
  ovvero \inlinecode{1}.

  \strategia{
    \begin{itemize}
      \item Facciamo la \textbf{ricorsione su} \inlinecode{n},
            ridotto di 1 ad ogni passo, con accumulatore \inlinecode{acc}
            che porta il prodotto parziale.
      \item \textbf{Ipotesi induttiva:} la funzione è corretta per qualsiasi
            valore di \inlinecode{n} strettamente minore di quello corrente.
      \item \textbf{Passo ricorsivo:} aggiorno \inlinecode{acc}
            moltiplicando per \inlinecode{n} e delego  nessun calcolo
            rimane pendente al ritorno.
            \vspace{0.1cm}

            \inlinecode{return fact\_helper(n - 1, acc * n);}

      \item \textbf{Caso base:} quando \inlinecode{n == 0} il calcolo
            è completo; viene restituito direttamente \inlinecode{acc}.
            \vspace{0.1cm}

            \inlinecode{if (n == 0) return acc;}
    \end{itemize}
  }

  \begin{codebox}
    \begin{lstlisting}
int fact_helper(int n, long long acc) {
    if (n == 0) return acc;
    return fact_helper(n - 1, acc * n);
}

int fact(int n) {
    return fact_helper(n, 1);   // acc inizializzato al valore neutro
}
    \end{lstlisting}
  \end{codebox}

\end{codeexample}

\subsection{ Accumulatore stringa}
Lo stato è una stringa che si costruisce progressivamente.


\begin{codeexample}[Inversione stringa con accumulatore]

  Per ottenere una tail recursion, si introduce un accumulatore
  \inlinecode{acc} di tipo \inlinecode{string} che costruisce la
  stringa invertita progressivamente, portandola in avanti invece
  di lasciarla pendente sullo stack. \inlinecode{acc} viene
  inizializzato alla stringa vuota, valore neutro della
  concatenazione.

  \strategia{
    \begin{itemize}
      \item Facciamo la \textbf{ricorsione sull'indice} \inlinecode{i},
            avanzato di 1 ad ogni passo, con accumulatore \inlinecode{acc}
            che porta la stringa invertita parziale.
      \item \textbf{Ipotesi induttiva:} la funzione è corretta per qualsiasi
            indice strettamente maggiore di \inlinecode{i}.
      \item \textbf{Passo ricorsivo:} prependo \inlinecode{s.at(i)}
            all'accumulatore e delego  nessun calcolo rimane
            pendente al ritorno.
            \vspace{0.1cm}

            \inlinecode{return inverti(s, i + 1, s.at(i) + acc);}

      \item \textbf{Caso base:} quando \inlinecode{i == (int)s.size()}
            tutti i caratteri sono stati processati; viene restituito
            direttamente \inlinecode{acc}.
            \vspace{0.1cm}

            \inlinecode{if (i == s.size()) return acc;}
    \end{itemize}
  }

  \begin{codebox}
    \begin{lstlisting}
string inverti(const string& s, int i = 0, string acc = "") {
    if (i == (int)s.size()) return acc;
    return inverti(s, i + 1, s.at(i) + acc);
}
    \end{lstlisting}
  \end{codebox}

\end{codeexample}
\subsection{ Accumulatori multipli paralleli}
Serve quando lo stato richiede più valori che evolvono insieme. I parametri accumulatori sono due o più e si aggiornano in parallelo ad ogni passo.

















\begin{codeexample}[Fibonacci con accumulatori paralleli]

  La successione di Fibonacci è definita come:
  \[
    F(n) = \begin{cases}
      0               & \text{se } n = 0 \\
      1               & \text{se } n = 1 \\
      F(n-1) + F(n-2) & \text{se } n > 1
    \end{cases}
  \]
  La definizione standard genera una ricorsione multipla con
  sottoproblemi sovrapposti. Per ottenere una tail recursion si
  introducono due accumulatori paralleli \inlinecode{a} e
  \inlinecode{b} che portano avanti gli ultimi due valori della
  sequenza, aggiornandosi ad ogni passo senza calcoli pendenti
  al ritorno.

  \strategia{
    \begin{itemize}
      \item Facciamo la \textbf{ricorsione su} \inlinecode{n},
            ridotto di 1 ad ogni passo, con due accumulatori
            paralleli \inlinecode{a} e \inlinecode{b} che portano
            rispettivamente $F(k)$ e $F(k+1)$ per il passo corrente.
      \item \textbf{Ipotesi induttiva:} la funzione è corretta per
            qualsiasi valore di \inlinecode{n} strettamente minore
            di quello corrente.
      \item \textbf{Passo ricorsivo:} aggiorno entrambi gli
            accumulatori in parallelo  \inlinecode{a} diventa
            \inlinecode{b} e \inlinecode{b} diventa \inlinecode{a+b}
             e delego. Nessun calcolo rimane pendente al ritorno.
            \vspace{0.1cm}

            \inlinecode{return fib(n - 1, b, a + b);}

      \item \textbf{Caso base:} quando \inlinecode{n == 0},
            \inlinecode{a} contiene già $F(0)$; viene restituito
            direttamente.
            \vspace{0.1cm}

            \inlinecode{if (n == 0) return a;}
    \end{itemize}
  }

  \begin{codebox}
    \begin{lstlisting}
int fib_helper(int n, long long a, long long b) {
    if (n == 0) return a;
    return fib_helper(n - 1, b, a + b);
}

    int fib(int n) {
    return fib_helper(n, 0, 1);   // a=F(0)=0, b=F(1)=1
}
    \end{lstlisting}
  \end{codebox}

\end{codeexample}













\section{Divide et impera}
\defn{}{
Divide et impera è una strategia ricorsiva che risolve un problema dividendolo in sotto problemi indipendenti della stessa natura, risolvendo ognuno per ipotesi induttiva,e poi combinando i risultati parziali in una soluzione completa.
}

Nel pattern Divide et impera ci sono due sottotipi, i quali sono:
\subsection{Divide et impera con combine esplicita}
Il lavoro principale avviene nella fase combina, dopo che entrambe la metà sono state risolte, si fa un lavoro non banale per unire i  risultati. 


\begin{codeexample}[Merge sort]








  Il \textbf{Merge Sort} è un algoritmo di ordinamento \textit{divide et impera}.
  Opera in due fasi:
  \begin{enumerate}
    \item \textbf{Divide:} trova il punto medio $m = \lfloor(l+r)/2\rfloor$ e
          richiama ricorsivamente \inlinecode{mergeSort} sulle due metà
          \inlinecode{[l,m]} e \inlinecode{[m+1,r]}.
    \item \textbf{Merge:} fonde le due metà ordinate confrontando elemento
          per elemento e copiando ogni volta il minore in un array di supporto,
          poi ricopia il risultato in \inlinecode{arr[l..r]}.
  \end{enumerate}


\begin{center}

\begin{tikzpicture}[
    nodes={minimum width=0.7cm, minimum height=0.7cm},
    row sep=-\pgflinewidth,
    column sep=-\pgflinewidth,
    >=stealth
]

% Livello 0: Array originale
\matrix(orig)[matrix of nodes,
    row 1/.style={nodes={draw=none}},
    nodes={draw}]
{
    \tiny{0} & \tiny{1} & \tiny{2} & \tiny{3}\\
    5 & 2 & 4 & 1\\
};

% Livello 1: Split
\matrix(left)[matrix of nodes,
    row 1/.style={nodes={draw=none}},
    nodes={draw},
    below left=1.2cm and 0.5cm of orig]
{
    \tiny{0} & \tiny{1}\\
    5 & 2\\
};

\matrix(right)[matrix of nodes,
    row 1/.style={nodes={draw=none}},
    nodes={draw},
    below right=1.2cm and 0.5cm of orig]
{
    \tiny{2} & \tiny{3}\\
    4 & 1\\
};

% Livello 2: Merge (ordinato)
\matrix(mleft)[matrix of nodes,
    row 1/.style={nodes={draw=none}},
    nodes={draw, fill=green!15},
    below left=1.2cm and 0.5cm of $(left)!0.5!(right)$]
{
    \tiny{0} & \tiny{1}\\
    2 & 5\\
};

\matrix(mright)[matrix of nodes,
    row 1/.style={nodes={draw=none}},
    nodes={draw, fill=green!15},
    below right=1.2cm and 0.5cm of $(left)!0.5!(right)$]
{
    \tiny{2} & \tiny{3}\\
    1 & 4\\
};

% Livello 3: Array finale ordinato
\matrix(final)[matrix of nodes,
    row 1/.style={nodes={draw=none}},
    nodes={draw, fill=orange!20},
    below=1.2cm of $(mleft)!0.5!(mright)$]
{
    \tiny{0} & \tiny{1} & \tiny{2} & \tiny{3}\\
    1 & 2 & 4 & 5\\
};

% Frecce split
\draw[->, thick, red!70!black] (orig-2-2.south) -- (left-1-1.north)
    node[midway, left, font=\scriptsize] {split};
\draw[->, thick, red!70!black] (orig-2-3.south) -- (right-1-2.north);

% Frecce merge
\draw[->, thick, green!50!black] (left-2-2.south) -- (mleft-1-1.north)
    node[midway, left, font=\scriptsize] {merge};
\draw[->, thick, green!50!black] (right-2-1.south) -- (mright-1-2.north);

% Frecce merge finale
\draw[->, thick, orange!70!black] (mleft-2-2.south) -- (final-1-2.north)
    node[midway, left, font=\scriptsize] {merge};
\draw[->, thick, orange!70!black] (mright-2-1.south) -- (final-1-3.north);

\end{tikzpicture}


\end{center}






    \strategia{
        \begin{itemize}
        \item Facciamo la \textbf{ricorsione sugli estremi} \inlinecode{l} e \inlinecode{r} dimezzando l'intervallo ad oggni passo. 
        \item \textbf{Ipotesi induttiva:}: la funzione è corretta per qualsiasi intervallo di dimensione strettamente minore di \inlinecode{r-l+1}
        \item \textbf{Passo ricorsivo:} per ipotesi le due metà sono già ordinate.
        \item \textbf{Caso base:} quando \inlinecode{l>=r} c'è un solo elemento ed è già ordinato
        \end{itemize}
    }



    \begin{codebox}
    \begin{lstlisting}
    void merge(vector<int>& v, int l, int m, int r) {
    vector<int> sx, dx;

    for (int i = l; i <= m; i++) sx.push_back(v.at(i));
    for (int i = m + 1; i <= r; i++) dx.push_back(v.at(i));

    int i = 0, j = 0, k = l;
    while (i < (int)sx.size() && j < (int)dx.size())
        v.at(k++) = (sx.at(i) <= dx.at(j)) ? sx.at(i++) : dx.at(j++);
    while (i < (int)sx.size()) v.at(k++) = sx.at(i++);
    while (j < (int)dx.size()) v.at(k++) = dx.at(j++);
}
    \end{lstlisting}
    \end{codebox}






    \begin{codebox}
    \begin{lstlisting}
    void mergeSort(vector<int>& v, int l, int r) {
    if (l >= r) return;
    int m = l + (r - l) / 2;     // DIVIDE
    mergeSort(v, l, m);           // CONQUER sinistra
    mergeSort(v, m + 1, r);       // CONQUER destra
    merge(v, l, m, r);            // COMBINE
}
    \end{lstlisting}
    \end{codebox}
\end{codeexample}


\begin{codeexample}[Massimo]
    \strategia{
    \begin{itemize}
    \item Facciamo la \textbf{riscorsione sugli estremi} \inlinecode{l} e \inlinecode{r}, dimmezando l'intervallo.
    \item \textbf{Ipotesi induttiva:} la funzione è coretta èer qualsiasi intervallo di dimensione strettamente minore di \inlinecode{r-l+1}.
    \item \textbf{Passo ricorsivo:} per ipotesi ho il massimo delle due metà; prendo il maggiore tra i due. 
    \item \textbf{Caso base:} quando \inlinecode{l==r} cìè un solo elemento 
    \end{itemize}
    }
    \begin{codebox}
    \begin{lstlisting}
    int massimo(const vector<int>& v, int l, int r){
    if(l==r) return v.at(l);
    int m = l+(r-1)/2;
    int maxSx = massimo(v, l, m);        // CONQUER sinistra
    int maxDx = massimo(v, m + 1, r);    // CONQUER destra
    return max(maxSx, maxDx);           // Combina 
    }
    \end{lstlisting}
    \end{codebox}
\end{codeexample}




\subsection{Divide et impera senza combina, con riduzione diretta}
Non c'è fase di combine  dopo aver diviso, si procede su una sola delle parti. È un caso degenere di Divide et impera che assomiglia alla ricorsione lineare, ma il criterio di divisione è $n/2$ invece di $ n-1$.

\begin{codeexample}[Fast power]
\strategia{
    \begin{itemize}
        \item Facciamo la \textbf{ricorsione} su \inlinecode{exp}, dimezzato ad ogni passo.
        \item \textbf{Ipotesi induttiva:} la funzione è coretta per qualsiasi valore di \inlinecode{exp} strettamente minore di quello corrente 
        \item \textbf{Passo ricorsivo:} per ipotesi ho \inlinecode{base\^(exp/2)}.
        \item \textbf{Caso base:}\ilinecode{exp== 0}
    \end{itemize}
}

    \begin{codebox}
    \begin{lstlisting}
    int fastPow(int base,int exp){
     if(exp== 0) return 1; 
     int half = fastPow(base,exp/2);
      if(exp%2 == 0) return half *hallf;
      else
            return base*half*half;
    }
    \end{lstlisting}
    \end{codebox}
\end{codeexample}

\begin{codeexample}
    Suppponiamo che qualcuno pensi a un numero tra 1 e 100, e noi dobbiamo indovinarlo. Ci verrà detto solo se è troppo alto , troppo basso, o se il numero è esatto.
    Per fare ciò ci sono due strategie:
    \begin{itemize}
    \item \textbf{Strategia  non ottimale}: Proviamo 1, poi 2, poi 3$\ldots$ se il numero fosse 99, ci vorebbero 99 tentativi. Ogni tentativo elimina un solo numero.
    \item \textbf{Strategia ottimale:} Ad ogni tentativo puntiamo sempre al centro dell'intervallo:
    \begin{enumerate}
        \item Proviamo $50 \rightarrow$ troppo alto elimino i numeri da 50 a 100.
        \item Proviamo $25 \rightarrow$ troppo alto elimino i numeri da 25 a 49.
        \item Proviamo $12 \rightarrow$ troppo basso limino i numeri da 1 a 12 .
        \item Proviamo $18 \rightarrow$ troppo basso limino i numeri da 13 a 18.
        \item Proviamo $21 \rightarrow$ troppo basso limino i numeri da 19 a 21.
        \item 23 trovato.
    \end{enumerate}
    \end{itemize}
    


    \begin{center}
        \begin{tikzpicture}[
    nodes={minimum width=1cm, minimum height=0.7cm},
    row sep=-\pgflinewidth,
    column sep=-\pgflinewidth,
    >=stealth
]

% === Passo 0: Array 1..100, mid=50 ===
\matrix(s0)[matrix of nodes, nodes={draw}]
{
    1 & \dots & |[fill=yellow!40]| 50 & \dots & 100\\
};
\node[right=0.5cm of s0, font=\scriptsize, align=left]
    {mid = 50\\ 23 < 50 $\Rightarrow$ vai a sx};

% === Passo 1: Array 1..49, mid=25 ===
\matrix(s1)[matrix of nodes, nodes={draw},
    below=1.2cm of s0]
{
    1 & \dots & |[fill=yellow!40]| 25 & \dots & |[fill=red!15]| 49\\
};
\node[right=0.5cm of s1, font=\scriptsize, align=left]
    {mid = 25\\ 23 < 25 $\Rightarrow$ vai a sx};

% === Passo 2: Array 1..24, mid=12 ===
\matrix(s2)[matrix of nodes, nodes={draw},
    below=1.2cm of s1]
{
    1 & \dots & |[fill=yellow!40]| 12 & \dots & |[fill=red!15]| 24\\
};
\node[right=0.5cm of s2, font=\scriptsize, align=left]
    {mid = 12\\ 23 > 12 $\Rightarrow$ vai a dx};

% === Passo 3: Array 13..24, mid=18 ===
\matrix(s3)[matrix of nodes, nodes={draw},
    below=1.2cm of s2]
{
    |[fill=red!15]| 13 & \dots & |[fill=yellow!40]| 18 & \dots & 24\\
};
\node[right=0.5cm of s3, font=\scriptsize, align=left]
    {mid = 18\\ 23 > 18 $\Rightarrow$ vai a dx};

% === Passo 4: Array 19..24, mid=21 ===
\matrix(s4)[matrix of nodes, nodes={draw},
    below=1.2cm of s3]
{
    |[fill=red!15]| 19 & \dots & |[fill=yellow!40]| 21 & \dots & 24\\
};
\node[right=0.5cm of s4, font=\scriptsize, align=left]
    {mid = 21\\ 23 > 21 $\Rightarrow$ vai a dx};

% === Passo 5: Array 22..24, mid=23 → TROVATO ===
\matrix(s5)[matrix of nodes, nodes={draw},
    below=1.2cm of s4]
{
    22 & |[fill=green!40]| 23 & 24\\
};
\node[right=0.5cm of s5, font=\scriptsize, align=left, text=green!50!black]
    {mid = 23\\ \textbf{Trovato!}};

% === Frecce tra i passi ===
\draw[->, thick] (s0.south) -- (s1.north);
\draw[->, thick] (s1.south) -- (s2.north);
\draw[->, thick] (s2.south) -- (s3.north);
\draw[->, thick] (s3.south) -- (s4.north);
\draw[->, thick] (s4.south) -- (s5.north);

% === Titolo e legenda ===
\node[above=0.4cm of s0, font=\bfseries]
    {Ricerca Binaria --- Cerco il 23 in [1, 100]};

\node[below=0.6cm of s5, font=\scriptsize, align=center]
    {
    \colorbox{yellow!40}{~mid~} = elemento centrale \quad
    \colorbox{red!15}{~scartato~} \quad
    \colorbox{green!40}{~trovato!~}\\[3pt]
    };

\end{tikzpicture}
    \end{center}
    \begin{codebox}
    \begin{lstlisting}
    bool ricercaBinaria(const vector<int>& v, int target, int da, int a){
        if(da>=a){
        int m = (da+a)/2
        if(v.at(m==2)
            return true;
        else 
            if(v.at(m)>e)
                return ricercaBinaria(v,e,da,m-1);
            else
                return ricercaBinaria(v,e,m+1,a);
        }
        else
            return false;
    }
    \end{lstlisting}
    \end{codebox}
\end{codeexample}

\section{Backtracking}
\defn{Backtracking}{
Il backtracking è una tecnica ricorsiva per esplorare sistematicamente uno spazio di scelte costruendo una soluzione incrementalmente. Ad ogni passo si effettua una scelta, si assume per ipotesi induttiva che la ricorsione esplori correttamente tutto ciò che segue, poi si annulla la scelta per esplorare le alternative.
}
Nel backtracking ci sono tre sottotipi i quali sono: 

\subsection{Backtracking su scelte binarie}
Ad ogni passo la scelta è binaria:includo l'elemento corrente oppure no. 

\begin{codeexample}[sottoinsiemi]
    \strategia{
        \begin{itemize}
        \item Facciamo la \textbf{ricorsione sull'iindice } \inlinecode{i}, avanzando ad ogni passo.
        \item \textbf{Passo ricorsivo:} er ogni elemento \inlinecode{v.at(i)} esploro i due rami — incluso e non incluso — usando la funzione corretta per ipotesi su \inlinecode{i+1}.
        \item \textbf{Caso base:} \inlinecode{i == v.size()} la soluzione pariziale è completa, e viene registrata.


        \end{itemize}
        }
    \begin{lstlisting}
    void subsets(const vector<int> & v, int i, vector<int> & corrente,vector<vector<int>> & res){
    if(i== v.size()){
        res.push_back(corrente);
        return;
    }
    subsets(v,i+1,corrente,res);
    corrente.push_back(v.at(i));
     subsets(v, i + 1, corrente, res);
      corrente.pop_back();   
    }
    \end{lstlisting}
    
\end{codeexample}

\subsection{ Backtracking su permutazioni}
Ad ogni passo le scelte sono filtrate da vincoli che dipendono dallo stato corrente.


\begin{codeexample}[Permutazioni]

  \strategia{
    \begin{itemize}
      \item Facciamo la \textbf{ricorsione sulla lunghezza} di
            \inlinecode{corrente}, aumentata di 1 ad ogni passo.
      \item \textbf{Ipotesi induttiva:} la funzione è corretta per
            qualsiasi stato con un elemento in più già scelto in
            \inlinecode{corrente}.
      \item \textbf{Passo ricorsivo:} per ogni elemento non ancora
            usato, lo aggiungo a \inlinecode{corrente}, uso la funzione
            (corretta per ipotesi) per esplorare il resto, poi annullo
            la scelta per provare le alternative.
            \vspace{0.1cm}

            \inlinecode{usato.at(i) = true;}
            \vspace{0.1cm}

            \inlinecode{corrente.push\_back(v.at(i));}
            \vspace{0.1cm}

            \inlinecode{permutazioni(v, usato, corrente, res);}
            \vspace{0.1cm}

            \inlinecode{corrente.pop\_back();}
            \vspace{0.1cm}

            \inlinecode{usato.at(i) = false;}

      \item \textbf{Caso base:} quando
            \inlinecode{(int)corrente.size() == (int)v.size()}
            la permutazione è completa; viene registrata.
            \vspace{0.1cm}

            \inlinecode{if ((int)corrente.size() == (int)v.size())}
            \inlinecode{res.push\_back(corrente);}
    \end{itemize}
  }

  \begin{codebox}
    \begin{lstlisting}
void permutazioni(const vector<int>& v, vector<bool>& usato,
                  vector<int>& corrente, vector<vector<int>>& res) {
    if ((int)corrente.size() == (int)v.size()) {
        res.push_back(corrente);
        return;
    }
    for (int i = 0; i < (int)v.size(); i++) {
        if (!usato.at(i)) {
            usato.at(i) = true;
            corrente.push_back(v.at(i));            
            permutazioni(v, usato, corrente, res);  
            corrente.pop_back();                    
            usato.at(i) = false;                  
        }
    }
}
    \end{lstlisting}
  \end{codebox}

\end{codeexample}







\subsection{ Backtracking con vincoli strutturali }
Ad ogni passo le scelte sono filtrate da vincoli che dipendono dallo stato corrente. 



\begin{codeexample}[Parentesi ben formate]

  \strategia{
    \begin{itemize}
      \item Facciamo la \textbf{ricorsione sulla lunghezza} di
            \inlinecode{corrente}, aumentata di 1 ad ogni passo,
            con contatori \inlinecode{open} e \inlinecode{close}
            che tracciano le parentesi già inserite.
      \item \textbf{Ipotesi induttiva:} la funzione è corretta per
            qualsiasi stato con un carattere in più già aggiunto a
            \inlinecode{corrente}.
      \item \textbf{Passo ricorsivo:} provo ad aggiungere
            \inlinecode{(} se \inlinecode{open < n}, e \inlinecode{)}
            se \inlinecode{close < open}; il pruning è implicito nelle
            condizioni. L'undo è implicito poiché \inlinecode{corrente}
            è passata per valore.
            \vspace{0.1cm}

            \inlinecode{if (open < n)}
            \vspace{0.1cm}

            \inlinecode{\ \ \ \ parentesi(n, open+1, close, corrente+"(", res);}
            \vspace{0.1cm}

            \inlinecode{if (close < open)}
            \vspace{0.1cm}

            \inlinecode{\ \ \ \ parentesi(n, open, close+1, corrente+")", res);}

      \item \textbf{Caso base:} quando
            \inlinecode{(int)corrente.size() == 2*n}
            la stringa è completa; viene registrata.
            \vspace{0.1cm}

            \inlinecode{if ((int)corrente.size() == 2*n)}
            \inlinecode{res.push\_back(corrente);}
    \end{itemize}
  }

  \begin{codebox}
    \begin{lstlisting}
void parentesi(int n, int open, int close,
               string corrente, vector<string>& res) {
    if ((int)corrente.size() == 2 * n) {
        res.push_back(corrente);
        return;
    }
    if (open < n)
        parentesi(n, open + 1, close, corrente + "(", res);
    if (close < open)
        parentesi(n, open, close + 1, corrente + ")", res);
}
    \end{lstlisting}
  \end{codebox}

\end{codeexample}
\chapter{Classi}



\intro{}



\section{Introduzione alle classi}
Una struct o classe permette di creare un tipo di dato personalizzato che raggruppa più informazioni insieme. Ad esempio, una \inlinecode{struct Prova} può contenere tre campi; un \inlinecode{float x}, un \inlinecode{int y} e una  \inlinecode{string s}. La cosa importante da capire è che la combinazione dei valori di tutti i campi in un dato momento rappresenta lo stato dell'oggetto. 
\begin{codebox}
    \begin{lstlisting}
    struct Punto{
        float x;
        int y;
        string s;
    }
    \end{lstlisting}
\end{codebox}
\subsection{Metodi}
Oltre ai campi, una \inlinecode{struct} può avere dei metodi, che sono semplicemente funzioni definite al suo interno. Ad esempio, \inlinecode{void stampa()} è un metodo che stampa \inlinecode{"hello world"} a schermo. La distinzione è netta: i campi descrivono cosa è l'oggetto, i metodi descrivono cosa sa fare.

Una volta definita la \inlinecode{struct}, si possono istanziare (creare) oggetti di quel tipo, ad esempio con Prova \inlinecode{o1, o2;}. Ogni oggetto è indipendente: \inlinecode{o1.x = 23.12} e \inlinecode{o2.x = 14.33} assegnano valori diversi allo stesso campo su istanze diverse. Per accedere a campi o metodi si usa l'operatore punto (.), come in \inlinecode{o1.stampa()}.




\subsection{Struttura generale di una classe}
Una classe può avere tre livelli di visibilità: \inlinecode{public}, \inlinecode{private} e \inlinecode{protected}.  Nella sezione \inlinecode{public} si mettono i metodi che l'utente della classe può usare dall'esterno. Nella sezione \inlinecode{private} si mettono sia i dati che eventuali metodi interni che non devono essere esposti. Quando un metodo viene definito direttamente dentrlo la classe non devono essere esposti. Quando un metodo viene definito direttamente dentro la classe si dice che è \textbf{inline}: dichiarazione e definizione coincidono nello stesso punto.

\subsection{Implementare i metodi fuori dalla classe}
Nelle classi i metodi si dichiarano dentro la classe e si definiscono separatamente usando l'operatore \inlinecode{::}. 

\subsection{Incapsulamento}
 I campi \inlinecode{double x} e \inlinecode{double y} vengono dichiarati private, il che significa che non possono essere modificati direttamente dall'esterno della \inlinecode{struct}. Per farlo, si devono usare i metodi \inlinecode{public} come \inlinecode{setx() }e \inlinecode{sety()}. In questo modo si nascondono i dati interni e si controlla come vengono modificati, evitando usi scorretti.

\begin{codebox}
    \begin{lstlisting}
    struct Punto{
        private:
            double x;
            double y;
        public:
        void setx(double nx);
        void sety(double ny);
        void setm(double nm);
        void setdg(double na);
        void  stampa();
    }
    \end{lstlisting}
\end{codebox}

\begin{center}
    \begin{figure}[H]
  \centering

\subsection{Coordinate}
Un punto nel piano può essere rappresentato in due modi matematicamente equivalenti:
\begin{itemize}
    \item \textbf{Cartesiano:} con $(x,y)$, le due coordinate lungo gli assi.
    \item \textbf{Polare:} con $(\rho,\theta)$, dove $\rho$ è la distanza dall'origine e theta è l'angolo rispetto all'asse $x$.
\end{itemize}
La \inlinecode{struct} Punto sfrutta proprio questa dualità: internamente salva sempre x e y come campi privati, ma espone anche i metodi \inlinecode{setm()} (per impostare il modulo $\rho$) e \inlinecode{setdg()}. Questi metodi eseguono internamente la conversione matematica:

\[
    x = \rho \cos \theta \quad y= \rho \sin \theta
\]

Chi usa la classe non deve sapere come i dati sono salvati internamente  può usare indifferentemente il sistema cartesiano o quello polare, e sarà a gestire la conversione.




  % ── Grafico 1: Coordinate Cartesiane ──────────────────────
  \begin{minipage}{0.45\textwidth}
    \centering
    \begin{tikzpicture}
      \draw[->] (-0.5, 0) -- (3.5, 0) node[right] {$x$};
      \draw[->] (0, -0.5) -- (0, 3.5) node[above] {$y$};

      \coordinate (P1) at (2.5, 2);

      \draw[dashed] (P1) -- (2.5, 0) node[below] {$x$};
      \draw[dashed] (P1) -- (0, 2)   node[left]  {$y$};

      \filldraw[black] (P1) circle (2pt);
      \node[below left] at (0,0) {$O$};
    \end{tikzpicture}
    \caption*{Coordinate cartesiane}
  \end{minipage}
  \hfill
  % ── Grafico 2: Coordinate Polari ──────────────────────────
  \begin{minipage}{0.45\textwidth}
    \centering
    \begin{tikzpicture}
      \draw[->] (-0.5, 0) -- (3.5, 0);
      \draw[->] (0, -0.5) -- (0, 3.5);

      \coordinate (P2) at (2.8, 2.2);

      \draw[-] (0,0) -- (P2) node[midway, above left] {$\rho$};
      \draw[->] (0.8, 0) arc[start angle=0, end angle=38, radius=0.8]
                node[right] {$\theta$};

      \filldraw[black] (P2) circle (2pt) node[above right] {$P$};
      \filldraw[black] (0,0) circle (1.5pt);
      \node[below left] at (0,0) {$O$};
    \end{tikzpicture}
    \caption*{Coordinate polari}
  \end{minipage}

  \caption{Rappresentazione di un punto nel piano}
\end{figure}
\end{center}


\subsection{Implementazione dei metodi della classe Punto}
Ora vedremo l'implementazione dei metodi della classe Punto:

\begin{codebox}
    \begin{lstlisting}
        void punto:: setx(double nx){
        x=nx;
        }
    \end{lstlisting}
\end{codebox}

\begin{codebox}
    \begin{lstlisting}
        void punto:: sety(double ny){
        y=ny;
        }
    \end{lstlisting}
\end{codebox}
\begin{codebox}
    \begin{lstlisting}
    void Punto::setrm(double nm) {
     double theta = atan2(y, x);     // recupera l'angolo corrente
     x = nm * cos(theta);
     y = nm * sin(theta);
}
    \end{lstlisting}
\end{codebox}


Il metodo \inlinecode{setm(double nm) }è un esempio di come l'incapsulamento protegga la coerenza interna dell'oggetto. L'obiettivo è cambiare il modulo $\rho$ mantenendo invariato l'angolo $\theta$: il punto si sposta lungo la stessa direzione ma a distanza diversa dall'origine. Senza incapsulamento, un utente potrebbe cambiare $x$ senza aggiornare $y$, rompendo la coerenza  con i metodi privati questo è impossibile.

\subsection{costruttore}

Il costruttore è un metodo  che viene chiamato automaticamente ogni volta che si crea un nuovo oggetto di una classe. Ha due caratteristiche che lo distinguono da tutti gli altri metodi: porta lo stesso nome della classe e non ha tipo di ritorno. Il suo scopo è garantire che l'oggetto nasca in uno stato valido, evitando che i campi contengano valori casuali presenti in memoria
\begin{codebox}
    \begin{lstlisting}
Punto::Punto() {
    x = 0;
    y = 0;
}\end{lstlisting}
\end{codebox}



\subsection{Distruttore}
Il distruttore è un metodo che viene chiamato automaticamente quando un oggetto viene eliminato, sia che esca dallo scope, sia che venga cancellato esplicitamente con \inlinecode{delete}. La sua sintassi è identica a quella del costruttore, ma con una tilde $\sim$ davanti al nome della classe, e anche lui non ha tipo di ritorno né parametri:
\begin{codebox}
    \begin{lstlisting}
Punto:: ~ Punto() {  
} \end{lstlisting}
\end{codebox}



\section{Live Coding  Classi in C++}

In questa lezione verrà svolta una sessione di live coding dedicata
all'introduzione delle classi in C++. Durante la lezione verranno
analizzati i concetti fondamentali legati alla definizione di una classe,
alla distinzione tra sezione \texttt{public} e \texttt{private}, e alla
gestione della memoria dinamica mediante costruttore, distruttore e
copy constructor.

L'obiettivo è mostrare in modo pratico il ciclo di vita di un oggetto,
la differenza tra allocazione su stack e su heap, e il motivo per cui
il passaggio per valore di oggetti che gestiscono memoria dinamica può
portare a errori come il \textit{segmentation fault}.

Il materiale completo della lezione, insieme al notebook e al codice
sviluppato durante la sessione, sarà disponibile nella repository del corso:

\noindent\textbf{Repository:} \url{https://github.com/ibrahimatelybarry23/pelmod1/tree/main/LiveCoding}

\end{document}
